% PLAN1-JP 2026-11-11 - verbatim Japanese source excerpts from 73-page main_JP.tex
% Source blob: 110f10f9342403d4c1c71dce9be2b7a4a5b9015a
% NOTE: fragments only, NOT a standalone PDF or included appendix.

% Section 7-8 excerpts include common core for reference; small-alphabet results are intact.
\section{ファンアウト1での境界}
\label{sec:fanout-one}

この節は定理~\ref{thm:C} を証明する．

\paragraph{三文字の倍増族．}
$\mathtt a,\mathtt b,\mathtt c$ を相異なる文字とし，$U_1=\{1,\zeta\}$ を，
単位元 $1$ と $\zeta^2=\zeta$ をもつものとし，
$s_U:\{\mathtt a,\mathtt b,\mathtt c\}^+\twoheadrightarrow U_1$ を，$\mathtt a,\mathtt b$
を $1$ に，$\mathtt c$ を $\zeta$ に送るものとする；したがって $s_U(w)=1$ となるのは，ちょうど
$\mathtt c$ を含まない空でない語である．$n\ge1$ について $G_n$ をファンアウト1の文法
\[
 S\to E_n,\qquad E_0\to\mathtt a,\qquad
 E_i\to\mathtt bE_{i-1}E_{i-1}\mathtt b\mid\mathtt cE_{i-1}\mathtt c\quad(1\le i\le n),
\]
とし，$L^{(i)}:=L_{E_i}(G_n)$ とおく．したがって $L^{(0)}=\{\mathtt a\}$ かつ
$L^{(i)}=\mathtt bL^{(i-1)}L^{(i-1)}\mathtt b\cup\mathtt cL^{(i-1)}\mathtt c$ である．また
$L_n:=L^{(n)}=L(G_n)$，$W_0:=\mathtt a$ および $W_i:=\mathtt bW_{i-1}W_{i-1}\mathtt b$ とおく．
これは
姉妹 CFG 原稿 \cite[Proposition~6.4]{kuriyamaCFG} の括弧付き倍増族の三文字の変種であり，そこでは
開き括弧と閉じ括弧が異なる文字である；その性質はすべて
ここで証明する．

\begin{proposition}[三文字の倍増族]
\label{prop:bracketed-family}
\label{prop:three-letter-family}
すべての $n\ge1$ について：
\begin{enumerate}[label=\textup{(\roman*)},leftmargin=*]
\item $G_n$ は既約かつ良く，ファンアウト1かつ規則ランク2であり，
$\|G_n\|=13n+7$ かつ $\bar\theta_{G_n}=4n$ をもつ．
\item $L_n$ は Clark--Eyraud の意味で代入可能である：すべての $x,y\in\Sigma^+$ について，
$\D^{(1)}_{L_n}(x)\cap\D^{(1)}_{L_n}(y)\ne\emptyset$ ならば
$\D^{(1)}_{L_n}(x)=\D^{(1)}_{L_n}(y)$ である．したがって $L_n\in\Cmcf{1}{s}$ が，すべての
型付け $s$（すべてのアルファベット $\Sigma\supseteq\{\mathtt a,\mathtt b,\mathtt c\}$ 上の）について成り立つ；
特に $L_n\in\Cmcf{1}{s_U}$．
\item $W_n$ は $L_n$ の語で $\mathtt c$ を含まない唯一のものであり，したがって
$s_U$-型 $1$ の唯一のものであり，$|W_n|=3\cdot2^n-2$．
\item すべての特徴標本 $F$（$L_n$ の）で，
$\mathcal B^{\le r}_{1,s_U}$ $(r\ge1)$ または $\mathcal B_{1,s_U}$ に対するもの，あるいは，
$\Cmcf{1}{s_U}$ 全体で特徴標本をもつ任意の集合駆動型作用素に対するものは，
$W_n$ を含む；したがって $\|F\|_+\ge3\cdot2^n-1$．
\end{enumerate}
\end{proposition}

証明は付録~\ref{app:bracketed-family} にある：(iii) が成り立つのは，
$\mathtt c$ を含まない導出が完全二分木であるからであり，(iv) は
(iii) と命題~\ref{prop:rare-type} から従う．


\begin{definition}[一様多項式特徴データ長]
\label{def:uniform-polydata}
$f,r_{\max}\ge1$ とする．有限半群の擬多様体 $\mathbf V$ が
\emph{ファンアウト $f$ かつ
ランク $r_{\max}$ で一様多項式特徴データ長}をもつとは，すべての空でない有限アルファベット $\Sigma$ と，すべての
全射な型付け $s:\Sigma^+\twoheadrightarrow\mathsf S$ で
$\mathsf S\in\mathbf V$ なるものについて，ある多項式 $p=p_{f,r_{\max},s,\Sigma}$ が存在して，
すべての $L\in\Cmcf{f}{s}$ と，すべての既約な良い
ファンアウト $\le f$，ランク $\le r_{\max}$ の表示 $G$（$L\cap\Sigma^+$ の）について，
作用素 $\mathcal B^{\le r_{\max}}_{f,s}$ が特徴標本
$F$ を $L$ に対して $\|F\|_+\le p(\|G\|,\bar\theta_G)$ なるものとしてもつことをいう．
\end{definition}

量化子は書かれたとおりの順序である：$p$ は $f$，
$r_{\max}$，$\Sigma$ および型付け $s$ に依存してよいが，標的 $L$ にも
表示 $G$ にも依存しない．

\begin{theorem}[ファンアウト1：データ長に対する擬多様体の境界]
\label{thm:fanout-one-data}
$r_{\max}\ge2$ とし，$\mathbf V$ を有限
半群の擬多様体とする．
\begin{enumerate}[label=\textup{(\roman*)},leftmargin=*]
\item $\mathbf V$ がファンアウト1かつランク $r_{\max}$ で一様多項式特徴データ長をもつのは，
$\mathbf V\subseteq\mathbf{LG}$ であるとき，かつそのときに限る．
\item $\mathbf V\subseteq\mathbf{LG}$ ならば，
定義~\ref{def:uniform-polydata} において
$p(\|G\|,\bar\theta_G)=O_{r_{\max},s,\Sigma}(\|G\|^2\bar\theta_G)$ ととってよい．
\item $\mathbf V\not\subseteq\mathbf{LG}$ ならば，
命題~\ref{prop:bracketed-family} の型付け $s_U$ は正の像
$U_1\in\mathbf V$ をもち，標的 $L_n\in\Cmcf{1}{s_U}$ は既約な良い
ランク2の表示 $G_n$ で $\|G_n\|=13n+7$ かつ
$\bar\theta_{G_n}=4n$ なるものをもつが，一方 $L_n$ のすべての特徴標本は
符号化総長少なくとも $3\cdot2^n-1\ge2^{(\|G_n\|+\bar\theta_{G_n})/17}$ をもつ．この下界は，
$\mathcal B^{\le r_{\max}}_{1,s_U}$ について，$\mathcal B_{1,s_U}$ について，および，
$\Cmcf{1}{s_U}$ 全体で特徴標本をもつすべての集合駆動型作用素について成り立つ．
\end{enumerate}
\end{theorem}

\begin{proof}
(ii)，したがって (i) の \emph{if} の方向は，
$f=1$ とした定理~\ref{thm:local-group-polydata} である．なぜならそのとき，すべての
$\mathsf S\in\mathbf V$ は $\mathbf{LG}$ に属するからである．(iii) については，
$\mathsf S\in\mathbf V\setminus\mathbf{LG}$ をとる．
定理~\ref{thm:lg-characterization}(ii) により，$\mathsf S$ は
$U_1$ と同型な部分半群をもち，擬多様体は部分半群について閉じているので
$U_1\in\mathbf V$ である．型付け $s_U$ は $U_1$ の上へ写す；
残りの主張は命題~\ref{prop:bracketed-family}(i) と~(iv) であり，
$\|G_n\|=O(n)$ と $\theta_{G_n}\le4n$ を用いる．
$3\cdot2^n-1$ は $\|G_n\|$ と
$\bar\theta_{G_n}$ のどんな多項式でも抑えられないので，(i) の \emph{only if} の方向が従う．
\end{proof}

\begin{remark}[定理~\ref{thm:fanout-one-data} の範囲]
\label{rem:fanout-one-data}
(a) 下界はランク2を用いる：倍増規則
$E_i\to\mathtt bE_{i-1}E_{i-1}\mathtt b$ は二分である．
線形 CFG で表される標的については，姉妹原稿は，すべての固定した型付けについて，
文法の大きさのみについて多項式の特徴データを与える
\cite[Theorem~8.3]{kuriyamaCFG}．
(b) この定理は擬多様体に関するものである．どの個々の型付けが
多項式のデータをもつかは，$U_1$ の上への型付け（系~\ref{cor:u1-classification}）を除き，ここでは特徴づけない．
(c) 上界は固定した作用素
$\mathcal B^{\le r_{\max}}_{1,s}$ に対するものであるが，下界は，クラス全体で特徴標本をもつ
すべての集合駆動型作用素について成り立つ．
\end{remark}

\paragraph{多項式時間・厚データ．}
Eyraud，Heinz および Yoshinaka \cite[Definition~2.12]{EyraudHeinzYoshinaka2016} は，
\cite{Yoshinaka2008} に従い，クラスが多項式時間・
厚データ同定可能（IPTtD）であるとは，あるアルゴリズムが，与えられた標本の大きさについて多項式の時間で，
それと無矛盾な仮説を返し，かつ，標的のすべての表現 $R$
が，大きさが $R$ の大きさとその厚さについて多項式の特徴標本をもつことをいう，と定める．
補題~\ref{lem:thickness-equivalence}(ii),(iii) により，彼らの
厚さ（非終端記号の最短の産出のうち最長のもの，$t_G$）と本稿の規則の
厚さは，$\|G\|$ を与えたとき多項式的に同値である．標本の大きさ
$\sum_{w\in K}|w|$ と $\|K\|_+$ は高々語の個数だけ異なる．彼らの
定義はアルゴリズムに，固定した表現のクラスにおける仮説を返すことを求める；
$\mathcal B^{\le r_{\max}}_{f,s}$ の生の出力は
規則 $\widehat S\to\eps$ を含み得るので，まずそれを取り除く．

\begin{lemma}[既約な良い仮説]
\label{lem:reduced-good-hypothesis}
$f$，$r_{\max}\ge1$ および $s$ を固定し，$K\subseteq\Sigma^*$ を有限で
$K\cap\Sigma^+\ne\emptyset$ なるものとし，$G^+_K$ を，
$\mathcal B^{\le r_{\max}}_{f,s}(K)$ から規則 $\widehat S\to\eps$ を除いたものとする．このとき
$G^+_K$ はファンアウト $\le f$ かつランク $\le r_{\max}$ の既約な良い MCFG であり，
$L(G^+_K)=L(\mathcal B^{\le r_{\max}}_{f,s}(K))\setminus\{\eps\}$ であり，$G^+_K$ は
$K$ から時間 $(1+\|K\|_+)^{O(fr_{\max})}$ で計算可能である．
\end{lemma}

\begin{proof}
Start と Link の規則は恒等の雛形をもち，Const の規則は
空でない語のタプルを書き，Comp の規則は良い雛形をもつ
（定義~\ref{def:composition-witness}）；したがって $G^+_K$ のすべての規則は良く，
どの非終端記号も空の成分をもつタプルを導出しない．したがって $\eps$ は
取り除いた規則によってのみ導出され，これが言語を与える．すべての $[\tuple x]$ は
その Const 規則により生産的であり，$\widehat S$ は $K$ が
空でない語をもつので生産的である．到達可能性については，$\tuple x\in\mathrm{Obs}_d(K)$ とし，例えば
$w=E[\tuple x]\in K$ で $E=v_0\blank_1v_1\cdots\blank_dv_d\in\mathcal E_d$ とする．このとき $w\ne\eps$ であり，
$(w)\in\mathrm{Obs}_1(K)$ が文脈 $\blank_1$ によって成り立ち，雛形
$\rho=(v_0\xi_1^1v_1\cdots\xi_1^dv_d)$ は良い：それは $\eps$-free であり，各
変数を一回ずつ順に用い，$v_1,\ldots,v_{d-1}$ が
空でないので非併合である．したがって $((w),\rho,(\tuple x))$ はランク1の合成証拠であり，
$\widehat S\to[(w)]\to\rho([\tuple x])$ は $[\tuple x]$ に到達する．ファンアウト，ランクおよび
実行時間は定理~\ref{thm:poly} のものである．
\end{proof}

\begin{corollary}[多項式時間・厚データ同定]
\label{cor:iptd}
$f\ge1$，$r_{\max}\ge1$ および型付け $s$ を固定する．$\mathcal R$ を
対 $(G,z)$ からなるものとする．ここで対は，ビット $z\in\{0,1\}$ と，既約な良いファンアウト $\le f$，
ランク $\le r_{\max}$ の MCFG $G$ または規則をもたない固定した文法 $G_\emptyset$ のいずれかとからなる；対は
$L(G)\cup\{\eps:z=1\}$ を表す．
$\Cmcf{f}{s}\cap\mathsf{MCFL}(f,r_{\max})$ の成員は
$(G,[\eps\in L])$ で表される．ここで $G$ は既約な良いファンアウト $\le f$，ランク $\le r_{\max}$ の
表示（$L\cap\Sigma^+$ の）であるか，または $G=G_\emptyset$ である，$L\cap\Sigma^+=\emptyset$ ならば．
\begin{enumerate}[label=\textup{(\roman*)},leftmargin=*]
\item $s(\Sigma^+)\in\mathbf{LG}$ ならば，
補題~\ref{lem:reduced-good-hypothesis} の $(G^+_K,[\eps\in K])$ を返し，または
$(G_\emptyset,[\eps\in K])$ を返す（$K\cap\Sigma^+=\emptyset$ ならば）アルゴリズムは，
$\Cmcf{f}{s}\cap\mathsf{MCFL}(f,r_{\max})$ を，クラス $\mathcal R$ について，多項式時間・厚データ同定する．
\item $f=1$ かつ $r_{\max}\ge2$ ならば，どのアルゴリズムも，その実行
時間がどうであれ，IPTtD のデータ条件を
$\Cmcf{1}{s_U}\cap\mathsf{MCFL}(1,r_{\max})$ について，クラス $\mathcal R$ に関して満たさない．
\item したがって，$r_{\max}\ge2$ と擬多様体 $\mathbf V$ について，
次は同値である：$\mathbf V\subseteq\mathbf{LG}$；すべての型付け $s$
で $s(\Sigma^+)\in\mathbf V$ なるものについて，クラス
$\Cmcf{1}{s}\cap\mathsf{MCFL}(1,r_{\max})$ はクラス $\mathcal R$ について IPTtD である．
\end{enumerate}
\end{corollary}

\begin{proof}
(i) 補題~\ref{lem:reduced-good-hypothesis} により，仮説は
$\mathcal R$ に属し，$\|K\|_+$ について多項式の時間で計算され，
$L(\mathcal B^{\le r_{\max}}_{f,s}(K))$ を表し，これは $K$ を含む（各 $w\in K\cap\Sigma^+$
はその Start と Const の規則によって導出される）．
定理~\ref{thm:local-group-polydata} は標本 $F$ で，長さ
$O(\|G\|^2\bar\theta_G)$ をもち，$L(\mathcal B^{\le r_{\max}}_{f,s}(K))=L$ が
$F\subseteq K\subseteq L$ であるときにはいつでも成り立つものを与えるので，そのとき仮説は $L$ を表す；また
$\bar\theta_G\le2\|G\|\max\{1,t_G\}$．
(ii) \cite[Definition~2.12]{EyraudHeinzYoshinaka2016} では，アルゴリズムは
標本を読み，標本は文字列の有限集合であるので，それは集合駆動型作用素である．そのような
アルゴリズムはクラスのすべての成員に対して特徴標本をもつ．なぜなら，すべての成員 $L$ で $L\cap\Sigma^+\ne\emptyset$ なるものは，既約な良い
表示でランク $\le r_{\max}$ のものをもつからである（命題~\ref{prop:good-normalization}）．
命題~\ref{prop:bracketed-family} の標的 $L_n$ はクラスに属し，
$L_n\setminus\{W_n\}$ も属する：それは
補題~\ref{lem:type-deletion} により $\Cmcf{1}{s_U}$ に属し，有限である．
命題~\ref{prop:rare-type} により，$L_n$ のすべての特徴標本は
$W_n$ を含み，その長さは $3\cdot2^n-2$ であるが，一方 $\|G_n\|=O(n)$ かつ $t_{G_n}\le\theta_{G_n}\le4n$．
(iii) は，$f=1$ とした (i) と (ii) から従う．なぜなら，$\mathbf{LG}$ に
含まれない擬多様体は $U_1$ を含むからである
（定理~\ref{thm:fanout-one-data}）．
\end{proof}

仮説はすべての標本の上で $\mathcal R$ に属する；
特徴標本を含まない標本の上では，その言語はクラスに属する必要はない．これは
\cite{Yoshinaka2011} の学習器と同様である．項目~(i) はすべてのファンアウトで成り立つ．
項目~(ii) はファンアウト1でのみ証明されている；ファンアウト2でのその類似は未解決問題~\ref{op:fanout-two} の一部であり，
族 $L_n$ はそれを証明できない（定理~\ref{thm:doubling-no-lift}）．

\paragraph{$U_1$ の上への個々の型付け．}
定理~\ref{thm:fanout-one-data} は擬多様体についての主張である．
ファンアウト1の個々の型付けについては，どの語が $U_1$ の写し $\{e,z\}$ を実現するかが問題になる：
文字が互いに素な型 $e$ の語は指数的なデータを強いる．
一方，型がそれ自身の冪だけで実現される孤立した文字（定義~\ref{def:isolated-letter}）は多項式のデータを与える（下の定理~\ref{thm:two-letter-fanout-one}）．
$U_1$ の上への型付けについては，この二つの場合で尽くされる（系~\ref{cor:u1-classification}）．

\begin{corollary}[ファンアウト1での型付けの水準の障害]
\label{cor:typing-level-obstruction}
$r_{\max}\ge2$ とし，$s:\Sigma^+\twoheadrightarrow\mathsf S$ を型付けとする．
相異なる冪等元 $e,z\in\mathsf S$ が $ez=ze=z$ をみたし，
語 $u,v,w\in\Sigma^+$ が $s(u)=s(v)=e$，$s(w)=z$ をみたし，\emph{文字が互いに素}である
（$uvw$ の中で同じ文字が二度現れない）とする．
$H_n$ を，三つの新しい文字の上の命題~\ref{prop:three-letter-family} の言語 $L_n$ の
$\mathtt a\mapsto u$，$\mathtt b\mapsto v$，$\mathtt c\mapsto w$ による像とし，
$G''_n$ を $G_n$ の像，すなわち規則 $S\to E_n$，$E_0\to u$ と
$E_i\to vE_{i-1}E_{i-1}v\mid wE_{i-1}w$ をもつ文法とする．
このとき $H_n\in\Cmcf{1}{s}$ であり，$G''_n$ は $H_n$ のランク2の reduced good な表示で
\[
 \|G''_n\|=(2|v|+2|w|+9)n+|u|+6,\qquad
 \bar\theta_{G''_n}\le2|u|+2|v|+4n|w|
\]
をみたし，$\mathcal B^{\le r_{\max}}_{1,s}$，$\mathcal B_{1,s}$，あるいは $\Cmcf{1}{s}$ 全体で特徴標本をもつ任意の集合駆動型の作用素についての
$H_n$ のすべての特徴標本の総符号長は $(|u|+2|v|)2^n-2|v|$ より大きい．
したがってこの一つの型付けについて，定義~\ref{def:uniform-polydata} のようにファンアウト1，ランク $r_{\max}$ での
特徴データの長さを抑える多項式は存在しない．
\end{corollary}

\begin{proof}
$L_n$ は有限で Clark--Eyraud 代入可能である（命題~\ref{prop:three-letter-family}）ので，
補題~\ref{lem:app-u1-code} から $H_n$ も Clark--Eyraud 代入可能である；
それは有限なので $H_n\in\Cmcf{1}{s}$ である．
文法は good であり，$E_i\Rightarrow^*w^iuw^i$ なので reduced である．
大きさは定義~\ref{def:standard-mcfg} から従う：開始規則の大きさは $3$，基底の規則は $1+|u|$，
各段の二つの規則は $2|v|+5$ と $2|w|+3$ である．
他の選択をすべて規則 $E_j\to wE_{j-1}w$ で補うと，各規則は長さ高々 $|u|+2n|w|$（開始規則と単項の規則），
あるいは $2|u|+2|v|+(2n+2i-4)|w|$（段 $i$ の二項の規則）の成功した使用をもつ．
$\{e,z\}$ では $e$ は単位元で $z$ は吸収元なので，$H_n$ の語が型 $e$ をもつのは，それが
$\mathtt c$ を含まない $L_n$ の語，すなわち $W_n$ の像であるとき，かつそのときに限る．
像 $h(W_n)$ の長さは $|h(W_i)|=2|h(W_{i-1})|+2|v|$ から $(|u|+2|v|)2^n-2|v|$ であり，
命題~\ref{prop:rare-type} はそれをそのようなすべての標本に入れることを強いる．
\end{proof}

$u,v,w$ が文字ならば，これは型 $e$ の二つの文字の場合である；
特に，型 $1$ の文字を少なくとも二つもつ $U_1$ の上へのすべての型付けに適用できる．全射性から型 $\zeta$ の文字があるからである．
型 $e$ の文字がないときにも適用できる：
単位元 $1$，$g^2=1$ と零元 $z$ をもつ $\mathsf S=\{1,g,z\}$ と，
型付け $\mathtt a,\mathtt b,\mathtt c,\mathtt d\mapsto g$，$\mathtt t\mapsto z$ について，
$u=\mathtt{ab}$，$v=\mathtt{cd}$，$w=\mathtt t$，$e=1$ ととればよい．
文字が互いに素であることなしに，型 $e$ の可換でない二つの語が常に十分かどうかは一般には未解決である；
二つの語が可換であるのはそれらが一つの語の冪であるとき，かつそのときに限るので，そのような語がないことは，
型 $e$ のすべての語が一つの原始的な語の冪であることを意味する．
$U_1$ の上への型付けではこの問いは生じない：型 $1$ の文字が二つあればこの系が適用でき，
一つ（例えば $\mathtt a$）ならば型 $1$ の語は $\mathtt a$ の冪であり，データは多項式である（系~\ref{cor:u1-classification}）．

$\Sigma=\{\mathtt a,\mathtt b\}$ とし，$s_{U,2}:\Sigma^+\twoheadrightarrow U_1$ を $\mathtt a$ を
$1$ に，$\mathtt b$ を $\zeta$ に写すもの（二文字の型付け）とする．
型 $1$ の文字が一つしかないこの型付けには系~\ref{cor:typing-level-obstruction} は適用できない；
それでも，そのデータは対数的な大きさの表示について $\bar\theta_G$ の2乗になりうる．

\begin{proposition}[ファンアウト1での二文字の型付けの2次のデータ]
\label{prop:binary-u1-quadratic}
$m,n\ge1$ について，$L_{m,n}\subseteq\{\mathtt a,\mathtt b\}^+$ を
$\{\mathtt a^{2^m},\mathtt b\}$ の語 $2^n$ 個の連接全体の集合とする．
\begin{enumerate}[label=\textup{(\roman*)},leftmargin=*]
\item $L_{m,n}\in\Cmcf{1}{s_{U,2}}$．
\item $L_{m,n}$ は規則のランクが2のファンアウト1の reduced good な表示 $G_{m,n}$ で
$\|G_{m,n}\|=6m+6n+9$，$\bar\theta_{G_{m,n}}=2^m+2^n-1$ をみたすものをもつ．
\item $\mathcal B_{1,s_{U,2}}$，$\mathcal B^{\le r}_{1,s_{U,2}}$（$r\ge1$），
および $\Cmcf{1}{s_{U,2}}$ 全体で特徴標本をもつすべての集合駆動型の作用素について，
$L_{m,n}$ のすべての特徴標本 $F$ は $\|F\|_+>2^{m+n}$ をみたす．
\end{enumerate}
したがって $m=n$ のとき $\|G_{n,n}\|=12n+9$，$\bar\theta_{G_{n,n}}=2^{n+1}-1$，
$\|F\|_+>\bar\theta_{G_{n,n}}^{\,2}/4$ である；特に，この型付けについてファンアウト1では，
定理~\ref{thm:local-group-polydata} と対照的に，$p$ を多項式とする $p(\|G\|)\,\bar\theta_G$ の形の評価は存在しない．
\end{proposition}

証明は付録~\ref{app:u1-typings} にある．純な語 $\mathtt a^{2^{m+n}}$ は塊の長さと塊の個数の積であり，
一方 $\bar\theta_G$ はそれらの和程度である．長い塊を使う最短の文は，他のすべての位置を $\mathtt b$ で埋めるからである．
$\mathtt a$ と $\mathtt b$ 以外の文字を含む語は $L_{m,n}$ の中で空の分布をもつので，
この命題は，$\mathtt a$ が型 $1$ の唯一の文字で $\mathtt b$ が型 $\zeta$ をもつ $U_1$ の上へのすべての型付けについて，そのまま成り立つ．
それでもファンアウト1でのこれらの型付けのデータは多項式であり（下の系~\ref{cor:two-letter-fanout-one}），
$s_{U,2}$ についてはすべてのファンアウト $f\ge2$ でもそうである（系~\ref{cor:two-letter-semilattice}）．

\paragraph{孤立した文字と二文字．}
定理~\ref{thm:fanout-one-data}(iii) の族は三文字を使い，そのうち二つは型 $1$ である．
型がそれ自身の冪だけで実現される文字は，どのアルファベットの上でも，すでに多項式のデータを与える．

\begin{definition}[孤立した文字]
\label{def:isolated-letter}
文字 $\mathtt a\in\Sigma$ が型付け $s$ について\emph{孤立している}とは，
$s^{-1}(s(\mathtt a))=\mathtt a^+$ であり，$\mathtt a$ を含み $\mathtt a^+$ に入らないすべての語が一つの共通の型をもつことである．
\end{definition}

例えば $\Sigma=\{\mathtt a,\mathtt b\}$ とし，$c$ を\emph{内容型付け}，すなわち $c(w)$ が $w$ の文字の集合であるような，
和集合についての三元の半束 $\{\{\mathtt a\},\{\mathtt b\},\Sigma\}$ の上への型付けとする；
$\mathtt a$ を含み $\mathtt a^+$ の外にある語の内容は $\Sigma$ なので，両方の文字が $c$ について孤立している．
$s:\Sigma^+\twoheadrightarrow U_1$ で $\mathtt a$ が型 $1$ の唯一の文字ならば，$s^{-1}(1)=\mathtt a^+$ であり，
他のすべての語は型 $\zeta$ をもつので，どのアルファベットの上でも $\mathtt a$ は孤立している．
三文字以上の上では，$\mathtt{ab}$ と $\mathtt{ac}$ の内容が異なるので，内容型付けについて孤立した文字はない．

\begin{theorem}[孤立した文字：ファンアウト1での短い純な代表元]
\label{thm:two-letter-fanout-one}
$\mathtt a$ を型付け $s:\Sigma^+\twoheadrightarrow\mathsf S$ について孤立した文字とし，
$L\in\Cmcf{1}{s}$，$G$ を $L\cap\Sigma^+$ のファンアウト1で任意のランクの reduced good な表示とする．
$G$ の，あるいはその二分化 $G'$（付録~\ref{app:two-letter-fanout-one}）の，$\mathtt a^+$ の生成語をもつすべての非終端記号 $A$ について，
型付き非終端記号 $A_{s(\mathtt a)}$ は
\[
 \tau_{A_{s(\mathtt a)}}\le C(\|G\|+2)^{22}\bar\theta_G^{\,8}
\]
をみたす．ここで $C<10^{43}$ は $\Sigma$，$s$，ランクによらない絶対定数である．
特に，$L$ が $\mathtt a^+$ の語を含めば，その最短のものの長さは高々 $C(\|G\|+2)^{22}\bar\theta_G^{\,8}$ である．
\end{theorem}

この定理は孤立した文字の型だけを評価する．他の型は，次の系の二つの場合について，補題~\ref{lem:app-h-other} で評価される．
また，$\mathtt a$ を含み $\mathtt a^+$ の外にあるすべての語が一つの型をもつという仮定は本質的である：
三文字以上の内容型付けではこれが成り立たず，データは指数的になりうる（下を見よ）．

\begin{corollary}[孤立した文字からの多項式のデータ]
\label{cor:two-letter-fanout-one}
$r_{\max}\ge2$ とし，$s$ を次のいずれかとする．
\begin{enumerate}[label=\textup{(\roman*)},leftmargin=*]
\item 正の像が半束である $\Sigma=\{\mathtt a,\mathtt b\}$ 上の型付け（例えば $c$ や $s_{U,2}$）．
\item 任意の有限アルファベットの上の型付け $s:\Sigma^+\twoheadrightarrow U_1$ で，型 $1$ の文字がちょうど一つのもの．
\end{enumerate}
$L\cap\Sigma^+$ の任意のランクのファンアウト1の reduced good な表示 $G$ をもつすべての $L\in\Cmcf{1}{s}$ は，
$\mathcal B^{\le r_{\max}}_{1,s}$ と $\mathcal B_{1,s}$ に共通の特徴標本 $K$ で
\[
 \|K\|_+\le C'r_{\max}(\|G\|+2)^{23}\bar\theta_G^{\,8}+\mathbf 1_{\{\eps\in L\}}
\]
をみたすものをもつ．ここで $C'<10^{45}$ は絶対定数である．
\end{corollary}

\begin{corollary}[ファンアウト1での $U_1$ の上への型付け]
\label{cor:u1-classification}
$r_{\max}\ge2$ とし，$s:\Sigma^+\twoheadrightarrow U_1$ を型付けとする．
定義~\ref{def:uniform-polydata} が各型付けに求める評価が，ファンアウト1，ランク $r_{\max}$ で $s$ について成り立つのは，
型 $1$ の文字がちょうど一つのとき，すなわちある文字 $\mathtt a$ について $s^{-1}(1)=\mathtt a^+$ のとき，かつそのときに限る．より詳しく：
\begin{enumerate}[label=\textup{(\roman*)},leftmargin=*]
\item 型 $1$ の文字がちょうど一つならば，データはすべてのランクの表示について系~\ref{cor:two-letter-fanout-one} のとおりに評価され，
命題~\ref{prop:binary-u1-quadratic} により，$p$ を多項式とするどの $p(\|G\|)\bar\theta_G$ でも評価されない；
\item 型 $1$ の文字が少なくとも二つあれば，系~\ref{cor:typing-level-obstruction} の標的は，
そのランク2の reduced good な表示の大きさと厚さについて指数的なデータを要する．
\end{enumerate}
\end{corollary}

\begin{proof}
$s$ は全射で $U_1$ は文字の型で生成されるので，型 $1$ の文字と型 $\zeta$ の文字がある．
(i) は，系~\ref{cor:two-letter-fanout-one}(ii) と，命題~\ref{prop:binary-u1-quadratic} およびその後の注意を，型 $\zeta$ の文字を $\mathtt b$ として使ったものである．
(ii) について，$\mathtt a\ne\mathtt b$ が型 $1$，$\mathtt c$ が型 $\zeta$ をもつとする；
冪等元 $e:=1$ と $z:=\zeta$ は相異なり $ez=ze=z$ をみたし，$u:=\mathtt a$，$v:=\mathtt b$，$w:=\mathtt c$ は
$s(u)=s(v)=e$，$s(w)=z$ をみたす文字が互いに素な語なので，系~\ref{cor:typing-level-obstruction} が適用できる．
\end{proof}

定理~\ref{thm:two-letter-fanout-one} と系~\ref{cor:two-letter-fanout-one} の証明は付録~\ref{app:two-letter-fanout-one} にある．
考え方は次のとおりである．非終端記号の最短の生成語 $\mathtt a^n$ の，節点の数が最小の導出 $T$ をとる．
一つのラベルの距離 $d$ の二つの出現から，純な塊の消去によって，$\mathtt a^d$ と可換で $z_0z_0\sim\mathtt a^{c_0}z_0$ をみたす，
$\mathtt a^+$ の外の基底 $z_0$ が得られる．
$T$ の重い道の上のラベルの $\mathtt a^+$ の外の短い生成語は，両側をそれぞれ高々 $d$ 文字で詰めると $\mathtt a^ez_0$ に潰れる；
孤立性が使われるのはちょうどここであり，それは $\mathtt a$ を含み $\mathtt a^+$ の外にあるすべての語を比較可能にする．
$\mathtt a^fz_0$ が与えられた分布の類に入る指数 $f$ の集合は，ある $p_0$ を法とする剰余類である；
$p_0$ が存在すれば $\mathtt a\sim\mathtt a^{1+p_0}$ となり，純な語は短い同値な代表元をもつ．
そうでなければ所属は符号の価値で決まる．各符号はそれまでの符号の和を超えるように選ばれ，これは符号一つあたり高々 $3$ 倍の費用で済む．
これらの符号の置換を，倍増族の場合と同様に，文法の大きさに対して数え上げる．

系~\ref{cor:two-letter-fanout-one} は，すべてのアルファベットについて量化する定理~\ref{thm:fanout-one-data}(i) に矛盾しない：
半束は $\mathbf{LG}$ に含まれない擬多様体をなし，失敗は三文字の上で起きる．
二文字上の半束の上へのすべての型付けの正の像は $c$ の像の商なので，系の (i) はそれらすべてを覆う．
三文字以上では，内容型付け $c_\Sigma$ について孤立が成り立たないことは必然である：
命題~\ref{prop:three-letter-family}(ii) から $\Sigma\supseteq\{\mathtt a,\mathtt b,\mathtt c\}$ について $L_n\in\Cmcf{1}{c_\Sigma}$ であり，
$W_n$ はその内容 $\{\mathtt a,\mathtt b\}$ の唯一の語なので，命題~\ref{prop:rare-type} は長さ少なくとも $3\cdot2^n-1$ のデータを強いる．
証明の中で失敗が起きるのは符号を潰すところである：一つの非終端記号の内容 $\{\mathtt a,\mathtt c\}$ の短い生成語と
内容 $\{\mathtt a,\mathtt b\}$ の生成語は同値である必要がない．
したがって，ファンアウト1では，$U_1$ の上への型付けが多項式のデータをもつのは型 $1$ の文字がちょうど一つのとき，かつそのときに限る（系~\ref{cor:u1-classification}）；
そのときデータは多項式であるが，$p(\|G\|)\,\bar\theta_G$ の形ではない（命題~\ref{prop:binary-u1-quadratic}）．

表~\ref{tab:scope} は，各主な主張について，固定するパラメータ，
作用素および付帯条件をまとめる．

\begin{table}[htbp]
\caption{主な主張の範囲．「標準」は
$\mathcal B^{\le r_{\max}}_{f,s}$ と $\mathcal B_{f,s}$ の両方（一つの共通の標本で）を意味する；
「クラス全体」は，クラス全体で特徴標本をもつ
すべての集合駆動型作用素を意味する．別に述べない限り，$G$ は $L\cap\Sigma^+$ の既約な良い表示であって，
ファンアウト $\le f$ かつランク $\le r_{\max}$ のものである．}
\label{tab:scope}
\centering
\footnotesize
\setlength{\tabcolsep}{3pt}
\renewcommand{\arraystretch}{1.15}
\begin{tabularx}{\linewidth}{@{}>{\raggedright\arraybackslash}p{0.2\linewidth}>{\raggedright\arraybackslash}p{0.22\linewidth}>{\raggedright\arraybackslash}p{0.17\linewidth}>{\raggedright\arraybackslash}X@{}}
\hline
主張 & 固定するもの & 作用素 & 評価と条件\\
\hline
定理~\ref{thm:A}（定理~\ref{thm:local-group-polydata}）
 & $f$，$r_{\max}$，$\Sigma$，$s$ で $s(\Sigma^+)\in\mathbf{LG}$ のもの
 & 標準
 & 上界 $O(\|G\|^2\bar\theta_G)$，すべての $L\in\Cmcf{f}{s}$ とすべての $G$ について\\
定理~\ref{thm:B}(i)（命題~\ref{prop:linear-factor-all-fanout}）
 & $f$，任意の $s$
 & 標準（$r\ge1$），クラス全体
 & 有限の標的に対する下界 $\Omega(\|G\|\bar\theta_G)$；$|\Sigma|\ge4f+1$（$|\Sigma|\ge2$，ただし $f=1$ の場合）を要する；
   ランク2の表示\\
定理~\ref{thm:B}(ii)（系~\ref{cor:trivial-exact}）
 & $f$，$r_{\max}\ge1$，$s_{\mathrm{triv}}$
 & $\mathcal B_{f,s_{\mathrm{triv}}}$
 & 上界 $O(\|G\|\bar\theta_G)$（Yoshinaka \cite{Yoshinaka2011}）；$r_{\max}\ge2$ かつ $|\Sigma|\ge4f+1$（$\ge2$，ただし $f=1$ の場合）のとき達成される\\
定理~\ref{thm:C}(i)（定理~\ref{thm:fanout-one-data}）
 & $f=1$，$r_{\max}\ge2$，一つの擬多様体
 & $\mathcal B^{\le r_{\max}}_{1,s}$（定義~\ref{def:uniform-polydata}）
 & 一様多項式データをもつのは $\mathbf V\subseteq\mathbf{LG}$ のとき，かつそのときに限る\\
定理~\ref{thm:C}(ii)
  & $f=1$，三文字上の $s_U$
 & 標準，クラス全体
 & 下界 $3\cdot2^n-1$，大きさと厚さが $O(n)$ の表示に対して\\
系~\ref{cor:typing-level-obstruction}
 & $f=1$，$r_{\max}\ge2$，文字が互いに素な $u,v,w$ をもつ型付け
 & 標準，クラス全体
 & 下界 $(|u|+2|v|)2^n-2|v|$，大きさと厚さが $O(n)$ の表示に対して\\
命題~\ref{prop:binary-u1-quadratic}
 & $f=1$，$s_{U,2}$，あるいは型 $1$ の文字が一つの $U_1$ の上への任意の $s$
 & 標準（$r\ge1$），クラス全体
 & 下界 $\bar\theta_G^{\,2}/4$，$\|G\|=O(\log\bar\theta_G)$\\
系~\ref{cor:two-letter-semilattice}（定理~\ref{thm:two-letter-content}）
 & $f\ge2$，$r_{\max}\ge2$，$\Sigma=\{\mathtt a,\mathtt b\}$，半束の上への $s$
 & 標準
 & 上界 $O_{f,r_{\max}}(\|G\|^3\bar\theta_G^{\,4}\log\bar\theta_G)$；Parikh 飽和でなければ $O_{f,r_{\max}}(\|G\|^2\bar\theta_G)$\\
命題~\ref{prop:sl-reduction}
 & $f$，$r_{\max}$
 & $\mathcal B^{\le r_{\max}}_{f,s}$
 & $\mathbf{LG}$ が $f$ で境界であるのは，ある $c_\Sigma$ が $f$ で一様な多項式のデータをもたないとき，かつそのときに限る\\
系~\ref{cor:two-letter-fanout-one}
 & $f=1$，$r_{\max}\ge2$；$\Sigma=\{\mathtt a,\mathtt b\}$ で半束の上への $s$，あるいは任意の $\Sigma$ で型 $1$ の文字が一つの $U_1$ の上への $s$
 & 標準
 & 上界 $O(r_{\max}(\|G\|+2)^{23}\bar\theta_G^{\,8})$（絶対定数），すべての $L\in\Cmcf{1}{s}$ と任意のランクのすべての $G$ について\\
系~\ref{cor:u1-classification}
 & $f=1$，$r_{\max}\ge2$，$U_1$ の上への $s$
 & 標準，クラス全体
 & 型 $1$ の文字がちょうど一つのとき，かつそのときに限り多項式のデータ\\
系~\ref{cor:iptd}
 & $f$，$r_{\max}$，$s$
 & 集合駆動型
 & $\mathbf{LG}$ について IPTtD，すべての $f$ で；$f=1$ ではちょうど $\mathbf{V}\subseteq\mathbf{LG}$ について\\
定理~\ref{thm:doubling-no-lift}
 & 任意の $f\ge2$，任意の $s$
 & なし（所属）
  & 定理~\ref{thm:C}(ii) の族とその文字が互いに素な像は $\Cmcf{f}{s}$ を離れ，その強制された語はそこで決してその型の唯一の語ではない；$f\ge2$ でのデータの評価はない\\
命題~\ref{prop:lg-needed-target}，\ref{prop:nonabelian-needed-target}
 & 字母と標的 $L_{\mathrm A}$，$L_{\mathrm B}$
 & なし（所属）
 & 非文脈自由，ファンアウト2，ランク2；$f\ge2$ で $\Cmcf{f}{s_{\mathrm A}}$，$\Cmcf{f}{s_{\mathrm B}}$ に入り，すべての $f\ge1$ でそれぞれどの窓--群型付け，どの窓--可換群型付けのクラスにも入らない；データの評価はない\\
未解決
 & 
 & 
 & $\mathbf{LG}$ が $f\ge2$ で境界であるかどうか（未解決問題~\ref{op:fanout-two}），すなわち $|\Sigma|\ge3$ のすべての $c_\Sigma$ が $f$ で一様な多項式のデータをもつかどうか\\
\hline
\end{tabularx}
\end{table}

\section{ファンアウト2：持ち上がらないことと二文字}
\label{sec:fanout-two}
\label{subsec:power-transfer}

この節は，定理~\ref{thm:C}(ii) の指数的な族と，
その文字が互いに素な準同型による像が，ファンアウト
2には持ち越されないことを示し（定理~\ref{thm:doubling-no-lift}），
未解決問題~\ref{op:fanout-two} を内容型付けに帰着させ（命題~\ref{prop:sl-reduction}），
すべてのファンアウト $f\ge2$ で二文字のアルファベットを決着させる（定理~\ref{thm:two-letter-content}）．

\begin{lemma}[二つのタプル成分にわたる冪の移送]
\label{lem:power-transfer}
$f\ge2$，$L\in\Cmcf{f}{s}$，$u\in\Sigma^+$ とし，
$i_0,p\ge1$ を $s(u)$ の最終的な指数と周期とする．すなわち
$s(u)^{t+p}=s(u)^t$ が $t\ge i_0$ について成り立つとする．
$i,j\ge i_0$，$\delta\in p\mathbb Z_{>0}$ であり，かつ $L$ のある語が
因子 $u^{i+\delta+j+1}$ を含むとする．このとき，すべての
$w_{\mathrm L},w_{\mathrm R}\in\Sigma^*$ と $w_{\mathrm M}\in\Sigma^+$ について，
\[
 w_{\mathrm L}u^iw_{\mathrm M}u^{\delta+j}w_{\mathrm R}\in L\quad\Longleftrightarrow\quad
 w_{\mathrm L}u^{i+\delta}w_{\mathrm M}u^jw_{\mathrm R}\in L.
\]
\end{lemma}
\begin{proof}
与えられた語を $yu^{i+\delta+j+1}z$ と書く．
タプル $\tuple x=(u^i,u^{\delta+j})$ と
$\tuple x'=(u^{i+\delta},u^j)$ は，最終的な周期性により，成分ごとに同じ $s$-型をもつ．
それらはまた，共通の許容な
受理文脈 $E=y\blank_1u\blank_2z$ をもつ．なぜなら
$E[\tuple x]=E[\tuple x']=yu^{i+\delta+j+1}z\in L$ だからである．
したがって $(f,s)$-条件はそれらの二成分の分布を等しくする．
それらを許容文脈 $w_{\mathrm L}\blank_1w_{\mathrm M}\blank_2w_{\mathrm R}$ で試す；
条件 $w_{\mathrm M}\ne\eps$ は本質的である．
\end{proof}

\begin{theorem}[倍増族は持ち上がらない]
\label{thm:doubling-no-lift}
$s:\Sigma^+\twoheadrightarrow\mathsf S$ を任意の固定した有限半群型付けとし，
$u_{\mathtt a},u_{\mathtt b},u_{\mathtt c}\in\Sigma^+$ を文字が互いに素なものとし，$\varphi$ を
準同型 $\mathtt a\mapsto u_{\mathtt a}$，$\mathtt b\mapsto u_{\mathtt b}$，
$\mathtt c\mapsto u_{\mathtt c}$ とし，$H_n:=\varphi(L_n)$ とする．ここで
$L_n$ は命題~\ref{prop:three-letter-family} の族である\textup(したがって $H_n=L_n$ である，
$\varphi$ が恒等写像のときは\textup)．$n\ge2|\mathsf S|+2$ かつ $f\ge2$ とする．
すべての $L\in\Cmcf{f}{s}$ で $\varphi(W_n)$ を含むものは，同じ
型，同じ Parikh ベクトルをもち，$u_{\mathtt c}$ の文字を含まない別の語を含む．特に
$H_n\notin\Cmcf{f}{s}$ であり，$\Cmcf{f}{s}$ のどの成員においても $\varphi(W_n)$ は
その型の唯一の語ではない．
\end{theorem}
\begin{proof}
準同型 $\varphi$ は単射であり，$H_n$ の語で
$u_{\mathtt c}$ の文字を含まないものは，$L_n$ の語で $\mathtt c$ を含まないものの像であるので，
$\varphi(W_n)$ はそのような唯一の語である．最終的な指数 $i_0\ge1$ と周期
$p\ge1$ を $s(u_{\mathtt b})$ について，$i_0+p\le|\mathsf S|+1$ となるように選ぶ；このとき $n\ge2i_0+p+1$．
$i=i_0$，$\delta=p$ および $j=n-i-\delta-1\ge i_0$ とおく．語
$\varphi(W_n)=u_{\mathtt b}\,\varphi(W_{n-1})\,\varphi(W_{n-1})\,u_{\mathtt b}$ は
$u_{\mathtt b}^nu_{\mathtt a}$ で始まるので，$u_{\mathtt b}^{i+\delta+j+1}$ を含む．
$u_{\mathtt b}^i$ を，最初のものの後の $i$ 個の $u_{\mathtt b}$ の複製とし，
$u_{\mathtt b}^{\delta+j}$ を，$\delta+j\le n-1$ 個の複製であって，最初の
$u_{\mathtt a}$（$\varphi(W_{n-1})$ の二番目の複製の）の直前にあるものとする．そしてそれの前には少なくとも
$n-1$ 個の $u_{\mathtt b}$ の複製がある．このとき
\[
 \varphi(W_n)=w_{\mathrm L}u_{\mathtt b}^i w_{\mathrm M}u_{\mathtt b}^{\delta+j}w_{\mathrm R}
\]
であり，ここで $w_{\mathrm L}=u_{\mathtt b}$ であり，$w_{\mathrm M}\ne\eps$ は最初の
$u_{\mathtt a}$ を含む．$L\in\Cmcf{f}{s}$ が $\varphi(W_n)$ を含むとする．
補題~\ref{lem:power-transfer} により，
$w':=w_{\mathrm L}u_{\mathtt b}^{i+\delta}w_{\mathrm M}u_{\mathtt b}^jw_{\mathrm R}\in L$．
二つの埋めたタプルは成分ごとに等しい型をもつので，$s(w')=s(\varphi(W_n))$；$w'$ は
$\varphi(W_n)$ の Parikh ベクトルをもち，$u_{\mathtt c}$ の文字を含まないが，
$u_{\mathtt b}^{n+\delta}u_{\mathtt a}$ で始まり，$u_{\mathtt b}^nu_{\mathtt a}$ で始まるのではない．したがって
$w'\ne\varphi(W_n)$ である．なぜなら $u_{\mathtt a}$ と $u_{\mathtt b}$ は共通の文字をもたないからである．
$L=H_n$ については，これは上で示した $\varphi(W_n)$ の一意性に反する．
\end{proof}

\begin{remark}
冪の移送は語の総長と Parikh ベクトルを保つので，
稀な型の代表が短いことは示さない；それは
命題~\ref{prop:three-letter-family} の族とその文字が互いに素な像を，
語を加えた後でも，排除するだけである．
\end{remark}

\paragraph{半束への帰着．}
$\mathbf{SL}$ を有限の半束（可換で冪等な半群）の擬多様体とする．
空でない有限のアルファベット $\Sigma$ について，\emph{内容型付け} $c_\Sigma$ は
$w\in\Sigma^+$ を $w$ に現れる文字の集合に写す．これは和集合についての $\Sigma$ の空でない部分集合の上への準同型である．

\begin{lemma}[粗い型付けは特徴標本を受け継ぐ]
\label{lem:typing-monotone}
$s\preceq t$ を型付け，$f,r\ge1$，$L\in\Cmcf{f}{s}$ とする．すべての有限の $K\subseteq L$ について
\[
 L\bigl(\mathcal B^{\le r}_{f,t}(K)\bigr)\subseteq
 L\bigl(\mathcal B^{\le r}_{f,s}(K)\bigr)\subseteq L
\]
であり，$\mathcal B_{f,t}$ と $\mathcal B_{f,s}$ についても同様である．
したがって $\mathcal B^{\le r}_{f,t}$（あるいは $\mathcal B_{f,t}$）についての $L$ のすべての特徴標本は
$\mathcal B^{\le r}_{f,s}$（あるいは $\mathcal B_{f,s}$）についての特徴標本である．
\end{lemma}

\begin{proof}
二つの文法は同じ非終端記号と同じ Start，Const，Comp 規則をもつ（定義~\ref{def:reconstruction}）．
$t$-文法の Link 規則は，共通の文脈をもち $t$-型が等しい観測されたタプルを結ぶ；
$\ker t\subseteq\ker s$ なのでそれらの $s$-型も等しく，それは $s$-文法の Link 規則である．
これが最初の包含を与える；二番目は命題~\ref{prop:reconstruction-soundness} である．
$F$ が $t$-作用素について特徴的で $F\subseteq K\subseteq L$ ならば，左辺は $L$ に等しく，したがって中辺もそうである．
\end{proof}

\begin{proposition}[ファンアウトの境界の半束への帰着]
\label{prop:sl-reduction}
$f\ge1$，$r_{\max}\ge1$ を固定し，一様な多項式の特徴データの長さをファンアウト $f$，ランク $r_{\max}$ で読む
（定義~\ref{def:uniform-polydata}）．
\begin{enumerate}[label=\textup{(\roman*)},leftmargin=*]
\item $\mathbf V\not\subseteq\mathbf{LG}$ であるすべての擬多様体は $\mathbf{SL}$ を含む．
したがって，一様な多項式の特徴データの長さをもつ擬多様体がちょうど $\mathbf{LG}$ に含まれるものであるのは，
$\mathbf{SL}$ がそれをもたないとき，かつそのときに限る．
\item $\mathbf{SL}$ が一様な多項式の特徴データの長さをもつのは，すべての空でない有限のアルファベット $\Sigma$ について
内容型付け $c_\Sigma$ が定義~\ref{def:uniform-polydata} の条件をみたすとき，かつそのときに限る．
\end{enumerate}
\end{proposition}

\begin{proof}
(i) 有限半群 $\mathsf S\notin\mathbf{LG}$ はアルファベット $\mathsf S$ の上のその評価型付けの正の像なので，
$U_1$ と同型な部分半群をもつ（定理~\ref{thm:lg-characterization}）；したがって $U_1\in\mathbf V$ である．
すべての有限の半束 $\mathsf H$ は $x\mapsto(\chi_z(x))_{z\in\mathsf H}$ によって $U_1^{\mathsf H}$ に埋め込まれる．
ここで $zx=z$ ならば $\chi_z(x)=1$，そうでなければ $\chi_z(x)=\zeta$ である
（$zxy=z$ であるのは $zx=z$ かつ $zy=z$ のとき，かつそのときに限るので準同型である）．したがって $\mathbf{SL}\subseteq\mathbf V$ である．
定義~\ref{def:uniform-polydata} の性質は，量化される型付けが少なくなるので，擬多様体からその部分擬多様体へ受け継がれる．
$\mathbf{SL}$ がそれをもたなければ，$\mathbf V\not\subseteq\mathbf{LG}$ であるすべての擬多様体もそれをもたない．
一方 $\mathbf V\subseteq\mathbf{LG}$ であるすべての擬多様体は定理~\ref{thm:local-group-polydata} からそれをもつ．
逆に，$U_1\notin\mathbf{LG}$ なので $\mathbf{SL}\not\subseteq\mathbf{LG}$ である．

(ii) $c_\Sigma$ の像は半束であり，一方の向きを与える．
他方の向きについて，$\mathsf S\in\mathbf{SL}$ とする $s:\Sigma^+\twoheadrightarrow\mathsf S$ をとる．
$\mathsf S$ は可換で冪等なので，$s(w)$ は $w$ に現れる文字 $a$ についての $s(a)$ の積であり，
$s\preceq c_\Sigma$，$\Cmcf{f}{s}\subseteq\Cmcf{f}{c_\Sigma}$ である（定義~\ref{def:class-cmcf} の後の項目~(a)）．
$L\in\Cmcf{f}{s}$ と表示 $G$ について，長さ高々 $p_{c_\Sigma}(\|G\|,\bar\theta_G)$ の
$\mathcal B^{\le r_{\max}}_{f,c_\Sigma}$ についての特徴標本は，補題~\ref{lem:typing-monotone} から
$\mathcal B^{\le r_{\max}}_{f,s}$ についての特徴標本である；$p_{s}:=p_{c_\Sigma}$ ととればよい．
\end{proof}

定理~\ref{thm:fanout-one-data} により，$r_{\max}\ge2$ のときファンアウト1で $\mathbf{SL}$ はこの性質をもたない．
ただし二文字の上では，$\mathbf{SL}$ のすべての型付けはそこで多項式のデータをもつ（系~\ref{cor:two-letter-fanout-one}）．
ファンアウト2ではこの問いは未解決であり（未解決問題~\ref{op:fanout-two}），
次の定理はすべてのファンアウト $f\ge2$ で二文字のアルファベットを決着させる．
稀な型の下界は内容型付けに移る：$s\preceq c_\Sigma$ について，$s$-型 $m$ に含まれるすべての内容の型は，
最短の語の長さが少なくとも $\mu_m(L)$（$s$ について計算したもの）である．

\paragraph{すべてのファンアウト $f\ge2$ での二文字．}
語の $\Sigma^+$ への所属がその Parikh ベクトルのみによるとき，$L$ は \emph{Parikh 飽和}であるという．

\begin{theorem}[すべてのファンアウト $f\ge2$ での二文字の内容型付け]
\label{thm:two-letter-content}
$\Sigma=\{\mathtt a,\mathtt b\}$，$c:=c_\Sigma$，$f\ge2$，$r_{\max}\ge2$ とし，
$D_f:=(6f+1)3^f$，
$C_{f,r_{\max}}:=(r_{\max}+1)\bigl(4r_{\max}(3f+1)+D_f\bigr)+3$，
$H_m:=\sum_{k=1}^m1/k\le1+\ln m$（$H_0:=0$）とおく．
すべての $L\in\Cmcf{f}{c}$ と，$L\cap\Sigma^+$ のファンアウト $\le f$，ランク $\le r_{\max}$ のすべての reduced good な表示 $G$ について，
$\tau(G,c)\le C_{f,r_{\max}}|V|^2\bar\theta_G^{\,4}H_{\bar\theta_G}$ であり，作用素
$\mathcal B^{\le r_{\max}}_{f,c}$ と $\mathcal B_{f,c}$ は共通の特徴標本 $K$ で
\[
 \|K\|_+\le4^{fr_{\max}}\bigl(3+r_{\max}C_{f,r_{\max}}\bigr)\|G\|^3\bar\theta_G^{\,4}H_{\bar\theta_G}
 +\mathbf 1_{\{\eps\in L\}}
\]
をみたすものをもつ．
$L$ が Parikh 飽和でなければ $t_{\mathrm{typ}}(G,c)\le D_f(r_{\max}+1)|V|\bar\theta_G$ であり，
$\|K\|_+\le4^{fr_{\max}}\bigl(3+r_{\max}D_f(r_{\max}+1)\bigr)\|G\|^2\bar\theta_G
+\mathbf 1_{\{\eps\in L\}}$ ととれる．
\end{theorem}

\begin{corollary}[二文字の上の半束型付け]
\label{cor:two-letter-semilattice}
$f\ge2$，$r_{\max}\ge2$ とし，$s:\{\mathtt a,\mathtt b\}^+\twoheadrightarrow\mathsf S$ を
$\mathsf S\in\mathbf{SL}$ である型付け，例えば $s_{U,2}$（$\mathtt a\mapsto1$，$\mathtt b\mapsto\zeta$）とする．
このとき，定義~\ref{def:uniform-polydata} が各型付けに求める評価が，この固定した型付けについて，
$\mathcal B^{\le r_{\max}}_{f,s}$ と $\mathcal B_{f,s}$ に対し，定理~\ref{thm:two-letter-content} の評価で，ファンアウト $f$，ランク $r_{\max}$ で成り立つ．
\end{corollary}

\begin{proof}
命題~\ref{prop:sl-reduction}(ii) の証明と同様に $s\preceq c$，$\Cmcf{f}{s}\subseteq\Cmcf{f}{c}$ である；
定理~\ref{thm:two-letter-content} と補題~\ref{lem:typing-monotone} を適用する．
\end{proof}

証明は付録~\ref{app:two-letter-content} にある．それは，すべてのファンアウト $f\ge2$ で成り立つアリティ1と2の条件を，
一つの文字の連を通して使う：ある文字 $x$ について $x^4$ が $L$ の語に現れれば，
所属は $x$-骨格（$x$ の各極大の連を一つの $x$ で置き換えたもの）と文字 $x$ の個数のみによる．
また，連 $xx$ の隣の他方の文字の二つの連は，Parikh 飽和を強いるか，禁じられたより長い連に合流する．
飽和していなければ，これは最短の型付き導出の重い節点の反鎖を抑える；
飽和していれば，与えられた Parikh ベクトルの語の節点数最小の導出を直接解析する．
この定理はファンアウトを分離しない：系~\ref{cor:two-letter-fanout-one} により，二文字の半束型付けはファンアウト1でも多項式のデータをもつので，
二文字の上ではそれらはすべてのファンアウトで多項式のデータをもつ．
三文字以上では，内容 $\{\mathtt a,\mathtt b\}$ の長い成分が長い連を含むとは限らず，問いは未解決である．

% Former Appendices H, I, J preserved here in full.
\section{ファンアウト1での孤立した文字}
\label{app:two-letter-fanout-one}

この付録は定理~\ref{thm:two-letter-fanout-one} と系~\ref{cor:two-letter-fanout-one} を証明する．
以下では $s:\Sigma^+\twoheadrightarrow\mathsf S$ は型付け，$\mathtt a\in\Sigma$ は $s$ について孤立した文字（定義~\ref{def:isolated-letter}），
$L\in\Cmcf{1}{s}$ とし，$\eps\notin L$ と仮定してよい．
$\Sigma^+$ の語は，$\mathtt a^+$ に入るとき\emph{純}，$\mathtt a$ を含まないとき\emph{$\mathtt a$ なし}，
それ以外，すなわち $\mathtt a$ と他の文字を含むとき\emph{$\mathtt a$ 混合}という．
孤立性は次の二つの事実を通してのみ使われる：
\begin{enumerate}[label=\textup{(I\arabic*)},leftmargin=*]
\item すべての純な語は型 $s(\mathtt a)$ をもち，この型の語は純な語である；
\item すべての $\mathtt a$ 混合の語は一つの型をもつ．
\end{enumerate}
$x\in\Sigma^+$ について $\D(x):=\D^{(1)}_L(x)$ と書き，$\D(x)=\D(y)$ のとき $x\sim y$ と書く．
これは $\Sigma^+$ 上の合同関係であり，分布が空の語どうしも $\sim$ で同値である．
したがって $x\equiv_{L,s}y$ であるのは，$x\sim y$ かつ $s(x)=s(y)$ のとき，かつそのときに限る．

\subsection{準備}

\begin{lemma}[消去]
\label{lem:app-h-cancel}
$x,y\in\Sigma^+$ は同じ型をもち，$z\in\Sigma^+$，$\D(zx)\ne\emptyset$，$zx\sim zy$ とする．
このとき $x\sim y$ である．右側についても同様である．
\end{lemma}

\begin{proof}
$(l,r)\in\D(zx)$ について $lzyr\in L$ となるので，$x$ と $y$ は文脈 $(lz,r)$ を共有する．
\end{proof}

この補題は，純な $x,y$ か $\mathtt a$ 混合の $x,y$ にだけ適用する．それらはそれぞれ (I1)，(I2) から同じ型をもつ．
非終端記号 $X$ が導出的な文脈 $(l,r)$ をもてば，$X$ の同じ型の二つの生成語は $(l,r)$ を共有し，$\sim$ で同値である（補題~\ref{lem:typed-homogeneity}）．
$X$ が実現する型 $t$ について，型付き非終端記号 $X_t$ の代表元の長さ，すなわち $X$ の型 $t$ の生成語と $\sim$ で同値な型 $t$ の語の最短の長さを $\tau(X,t)$ と書く．

\paragraph{二分化．}
$G$ を $L$ のファンアウト1の reduced good な表示とし，$\sigma_G$ をその規則に現れる文字の個数とする．したがって $\sigma_G\le\min\{|\Sigma|,\|G\|\}$ である．
$t\ge3$ のすべての規則 $X\to Y_1\cdots Y_t$ を新しい非終端記号を通る右線形の鎖に置き換え，
長さ2以上の右辺に現れる終端記号を前終端記号 $T_{\mathtt x}\to\mathtt x$（使われる文字 $\mathtt x$ についてだけ）に置き換え，trim する．
得られる文法 $G'$ は reduced かつ good であり，規則は $X\to YZ$，$X\to Y$，$X\to\mathtt x$（$\mathtt x\in\Sigma$）の形で，$L$ を生成する．
$G$ の規則 $X\to\gamma$ は高々 $\max\{0,|\gamma|-2\}$ 個の新しい非終端記号と高々 $\max\{1,|\gamma|-1\}$ 個の規則を与え，
どちらも $\|G\|$ でのその規則の大きさ $1+r+|\gamma|$ 以下なので
\[
 k_V:=|V'|\le\|G\|+\sigma_G\le2\|G\|,\qquad |P'|\le\|G\|+\sigma_G
\]
であり，$|\Sigma|=2$ ならば $k_V\le\|G\|+2$，$|P'|\le\|G\|+2$ である．
さらに $\theta:=\bar\theta_{G'}\le\bar\theta_G$ である．新しい規則を使う導出は，それが由来する $G$ の規則を使う導出に対応し，
前終端の規則 $T_{\mathtt x}\to\mathtt x$ を使う導出は，$\mathtt x$ を書く $G$ の規則を使う導出に対応するからである．
$G'$ のすべての非終端記号は，長さ $\le\theta$ の生成語と，$|l|+|r|\le\theta$ の導出的な文脈 $(l,r)$ をもつ：
その非終端記号を左辺に，あるいは右辺にもつ規則を通る最短の文をとればよい．
$G$ の非終端記号は，同じ生成語をもつ $G'$ の非終端記号である．

\subsection{木と基底}

純な生成語をもつ $G'$ の非終端記号 $A$ を固定する．
その最短のものを $\mathtt a^n$，$A$ からの $\mathtt a^n$ の導出木で節点の数が最小のものを $T$，$|l|+|r|\le\theta$ をみたす $A$ の導出的な文脈を $(l,r)$ とする．
$T$ の部分木を同じラベルからの導出で置き換えると $A$ の生成語 $w$ が得られ，したがって $lwr\in L$ である．
生成語が $\theta$ より長い節点を\emph{重い}といい，重い子をもたない重い節点を\emph{種}という．

\begin{lemma}
\label{lem:app-h-tree}
\begin{enumerate}[label=\textup{(\alph*)},leftmargin=*]
\item $T$ の根から葉への道でラベルは繰り返さない；そのような道の非終端記号の節点は高々 $k_V$ 個である．
\item $T$ の同じラベル $Y$ の節点はすべて同じ生成語 $\mathtt a^{n_Y}$ をもつ．
\item 重い節点のラベル $Y$ はすべて，長さ $\le\theta$ の生成語 $\omega_Y\notin\mathtt a^+$ をもつ．
\item 種の生成語の長さは $\le2\theta$ である；種どうしは互いに素である；
$k\ge1$ 個の種を含む部分木の生成語の長さは $\le k(k_V+2)\theta$ である；
隣り合う二つの種の左端の差は高々 $2k_V\theta$ である．
\end{enumerate}
\end{lemma}

\begin{proof}
(a) 祖先と子孫が同じラベルをもてば，祖先の部分木を子孫の部分木に置き換えると，$A$ のより短い純な生成語か，より小さい木が得られる．
(b) そうでなければ長い方の部分木を短い方に置き換える．
(c) ラベル $Y$ の重い節点での規則を $\rho$ とする．
$\rho$ を使う長さ $\le\theta$ の文は，$|y|\le\theta<n_Y$ をみたす $Y$ の生成語 $y$ を含む；
$y\in\mathtt a^+$ ならば，それを $T$ に挿入すると $A$ のより短い純な生成語が得られる．
(d) 種は単位規則の節点ではなく，二つの軽い子をもつ．
種 $v$ が種 $u$ の下にあれば，$v$ の上にある $u$ の子は重い．
$k$ 個の種をもつ部分木の重い節点は種への道の上にあるので高々 $kk_V$ 個であり，種でないものはそれぞれ高々一つの軽い子をもつ．
隣り合う種 $u,v$ の間には，それらの最小共通祖先からの二本の道からぶら下がる部分木の生成語があり，それは高々 $2(k_V-1)$ 個である；
重いものがあれば $u$ と $v$ の間に種を含むので，どれも軽い．
\end{proof}

$X_0$ を $T$ の根とし，$X_{i+1}$ を $X_i$ の子のうち生成語の長い方とする．すると $n_{X_i}\le2n_{X_{i+1}}$ である．
根の部分木が種をもたなければ根は軽く $n\le\theta$ である；種を少なくとも一つ，高々 $k_V$ 個もてば $n\le k_V(k_V+2)\theta$ である．
そうでなければ，部分木が少なくとも $k_V+1$ 個の種を含む最大の添字を $i^*$ とする；
このとき $X_{i^*+1}$ が存在し，補題~\ref{lem:app-h-tree}(a) から $i^*\le k_V-2$ であり，$B_0:=n_{X_{i^*}}\le2k_V(k_V+2)\theta$ である；
$n_{X_{i^*}}<n$ と仮定してよく，$T$ は二項の節点をもつので $k_V\ge2$ である．
$X_{i^*}$ の下の種は $T$ の中で連続している；連続する $k_V+1$ 個のうち二つは同じラベル $Y$ をもち，
位置 $x_1<x_2$ にあって $d:=x_2-x_1\le2k_V^2\theta$ である．
$g_0:=d-n_Y\ge0$，$\omega:=\omega_Y$，$x:=x_1$，$y:=n-x_2-n_Y$，$y_1:=n-x_1-n_Y=d+y$ とおく．
二つの種の両方，最初のもの，二番目のものを $\omega$ で置き換えると，$A$ の生成語
\[
 W_1=\mathtt a^x\omega\mathtt a^{g_0}\omega\mathtt a^y,\qquad
 W_2=\mathtt a^x\omega\mathtt a^{d+y},\qquad
 W_3=\mathtt a^{x+d}\omega\mathtt a^y
\]
が得られる．ここで $x+g_0+y\ge1$ である：そうでなければ二つの種が，少なくとも $k_V+1\ge3$ 個の種を含む根の生成語を覆ってしまう．
$d\ge n_Y\ge1$ かつ $\omega\notin\mathtt a^+$ なので，語 $W_1,W_2,W_3$ は $\mathtt a$ 混合である．

\begin{lemma}[基底]
\label{lem:app-h-base}
$g_0\ge1$ または $\omega$ が $\mathtt a$ 混合ならば $z_0:=\omega\mathtt a^{g_0}$，$c_0:=d$ とおき，
$g_0=0$ かつ $\omega$ が $\mathtt a$ なしならば $z_0:=\omega\mathtt a^d$，$c_0:=2d$ とおく．
このとき $z_0$ は $\mathtt a$ 混合で，$\D(z_0)\ne\emptyset$，$d\mid c_0$ であり，
\[
 \mathtt a^dz_0\sim z_0\mathtt a^d\qquad\text{かつ}\qquad z_0z_0\sim\mathtt a^{c_0}z_0.
\]
\end{lemma}

\begin{proof}
(I2) から $W_1,W_2,W_3$ は一つの型をもつ；それらは $(l,r)$ を共有するので $\sim$ で同値である．
$W_2\sim W_3$ の左の $\mathtt a^x$ と右の $\mathtt a^y$ を消去すると（補題~\ref{lem:app-h-cancel}；残る語は $\mathtt a^d$ と $\omega$ を含むので $\mathtt a$ 混合であり，
消去の前の語は $W_2$ の因子なので分布は空でない），$\omega\mathtt a^d\sim\mathtt a^d\omega$ を得る．
$g_0\ge1$ または $\omega$ が $\mathtt a$ 混合ならば，$W_1\sim W_2$ で同じ消去を行うと（このとき $\omega\mathtt a^{g_0}\omega$ は $\mathtt a$ 混合），
$\omega\mathtt a^{g_0}\omega\sim\omega\mathtt a^d$ となり，$z_0z_0\sim\omega\mathtt a^{d+g_0}=z_0\mathtt a^d$ と $\mathtt a^dz_0\sim z_0\mathtt a^d$ が従う．
$g_0=0$ かつ $\omega$ が $\mathtt a$ なしならば，$n_Y=d$ かつ $x+y\ge1$ である．
$y\ge1$ ならば，$\mathtt a^x$ と $\mathtt a^{y-1}$ を消去して $\omega\omega\mathtt a\sim\omega\mathtt a^{d+1}$ を得るので，
$z_0z_0\sim\omega\omega\mathtt a^{2d}\sim\omega\mathtt a^{3d}=z_0\mathtt a^{2d}$ である；
$y=0$ ならば，$\mathtt a^{x-1}$ を消去して $\mathtt a\omega\omega\sim\mathtt a\omega\mathtt a^d$ を得るので，
$\omega\omega$ の前に文字 $\mathtt a$ を一つ残して
$z_0z_0\sim\mathtt a^{2d}\omega\omega=\mathtt a^{2d-1}(\mathtt a\omega\omega)\sim\mathtt a^{2d-1}(\mathtt a\omega\mathtt a^d)=\mathtt a^{2d}z_0$ である；
$\mathtt a$ なしの語 $\omega\omega$ は何とも比較されない．
ここで比較する語は $\mathtt a$ と $\omega$ を含むので $\mathtt a$ 混合であり，消去の前の語は $W_1$ の因子である．
どちらの場合も $z_0\mathtt a^{2d}\sim\mathtt a^{2d}z_0$ かつ $\mathtt a^dz_0\sim z_0\mathtt a^d$ である．
最後に，$z_0$ は $W_1$ または $W_2$ の因子である．
\end{proof}

$\mathtt a$ が型 $1$ の唯一の文字であるような $U_1$ の上への型付けでは，$\mathtt a$ なしの語は $\mathtt a$ 混合の語と同じ型をもち，
二番目の場合は避けられる；これは使わない．

\emph{注釈付きの符号}とは，$e\ge0$，$d\mid e$，$z\sim\mathtt a^ez_0$ をみたす組 $(z,e)$ であり，その\emph{価値}は $e+c_0$ である．
$d\mid R$ をみたす純な塊 $\mathtt a^R$ の価値は $R$ とする．
注釈付きの符号とそのような塊の指定された積 $u$ について，$\operatorname{val}(u)$ はその因子の価値の和である．
（一般に指数 $e$ は $z$ から定まらないので，価値は注釈に付ける．）

\begin{lemma}
\label{lem:app-h-collapse}
$u$ が少なくとも一つの符号を含めば，$u\sim\mathtt a^{\operatorname{val}(u)-c_0}z_0$ である．
\end{lemma}

\begin{proof}
各符号を $\mathtt a^ez_0$ で置き換え，$\mathtt a^dz_0\sim z_0\mathtt a^d$ ですべての塊を左へ移し，
$z_0z_0\sim\mathtt a^{c_0}z_0$ と $d\mid c_0$ から帰納法で従う $z_0^t\sim\mathtt a^{(t-1)c_0}z_0$ を使う．
\end{proof}

\subsection{重い道に沿った符号}

$n_{X_i}<n$ となる最小の添字を $i_0$ とする；$1\le i_0\le i^*$ である．
$i_0\le i\le i^*$ について，$T$ の中での $X_i$ の文脈を $(\mathtt a^{p_i},\mathtt a^{q_i})$ とすると $p_i+q_i\ge1$ である；
$X_i$ の部分木は二つの種を含む；その中で最初の種の左と右にある文字の数を $L_i,R_i$ とすると $R_i\ge g_0+n_Y\ge1$ であり，
$X_i$ の $\mathtt a$ 混合の生成語 $\Omega_i:=\mathtt a^{L_i}\omega\mathtt a^{R_i}$ をとる．

\begin{lemma}
\label{lem:app-h-borrow}
$\omega_{X_i}$ が $\mathtt a$ 混合ならば $\omega_{X_i}\sim\Omega_i$ である．
$\omega_{X_i}$ が $\mathtt a$ なしで $p_i\ge1$ ならば $\mathtt a\omega_{X_i}\sim\mathtt a\Omega_i$ である；
$\omega_{X_i}$ が $\mathtt a$ なしで $p_i=0$ ならば $\omega_{X_i}\mathtt a\sim\Omega_i\mathtt a$ である．
\end{lemma}

\begin{proof}
$l\mathtt a^{p_i}\omega_{X_i}\mathtt a^{q_i}r$ と $l\mathtt a^{p_i}\Omega_i\mathtt a^{q_i}r$ はともに $L$ に入る．
一番目の場合，$\omega_{X_i}$ と $\Omega_i$ は $\mathtt a$ 混合で文脈 $(l\mathtt a^{p_i},\mathtt a^{q_i}r)$ を共有する．
二番目の場合，$\mathtt a\omega_{X_i}$ と $\mathtt a\Omega_i$ は $\mathtt a$ 混合で文脈 $(l\mathtt a^{p_i-1},\mathtt a^{q_i}r)$ を共有する；
三番目の場合は $q_i\ge1$ として対称である．(I2) で結論する．
\end{proof}

$s'_i,s_i\in[0,d]$ を次のようにとる：$d\mid L_i+s'_i$；
補題~\ref{lem:app-h-base} の第一の場合は $R_i+s_i\equiv g_0\pmod d$，第二の場合は $R_i+s_i\equiv0\pmod d$；
補題~\ref{lem:app-h-borrow} で一文字借りるときは，それぞれ $s'_i\ge1$，$s_i\ge1$．
$\hat z_i:=\mathtt a^{s'_i}\omega_{X_i}\mathtt a^{s_i}$ とおく．
補題~\ref{lem:app-h-borrow} から $\hat z_i\sim\mathtt a^{s'_i+L_i}\omega\mathtt a^{R_i+s_i}$ である．
$R_i\ge g_0+n_Y$ なので，第一の場合は $k\ge0$ として $R_i+s_i=g_0+kd$，第二の場合は $k\ge0$ として $R_i+s_i=d+kd$ である；
したがって $e_i:=s'_i+L_i+kd$ について $\hat z_i\sim\mathtt a^{e_i}z_0$ であり，$(\hat z_i,e_i)$ は注釈付きの符号で，
どちらの場合も（$n_{X_i}=L_i+n_Y+R_i$ と，それぞれ $d-g_0=n_Y$，$n_Y=d$ を使って）その価値は
\[
 v_i=n_{X_i}+s_i+s'_i,\qquad|\hat z_i|\le\theta+2d
\]
である．$i<i_0$ については $v_i:=n$ とおく（符号はない）．
$X_i$ は二つの種を含む生成語をもつ $X_{i^*}$ の祖先なので，$i\le i^*$ について $n_{X_i}\ge d+n_Y$ である．
したがって，すべての $i\le i^*$ について
\begin{equation}
\label{eq:app-h-star}
 v_i\ge n_{X_i}>d,\qquad v_i\le2v_{i+1}+2d\ \ (i<i^*)
\end{equation}
であり，$v_{i^*}\le B_0+2d$ である；さらに $n\le2n_{X_{i_0}}\le2v_{i_0}$ である．

$A$ の生成語 $\mathtt a^{x_1}\omega\mathtt a^{y_1}$ は，$h:=y_1-g_0$ あるいは $h:=y_1-d$ として $\mathtt a^{x_1}z_0\mathtt a^{h}$ に等しい．
$0\le\rho,\rho'<d$，$d\mid e^*$ として $h=qd+\rho$，$x_1+qd=e^*+\rho'$ と書くと，
$l_\circ:=l\mathtt a^{\rho'}$，$r_\circ:=\mathtt a^\rho r$ について $l_\circ\mathtt a^{e^*}z_0r_\circ\in L$ を得，
$V^*:=e^*+c_0=n-\rho-\rho'$ である．

\subsection{周期性と所属}

\begin{lemma}
\label{lem:app-h-period}
$P_0:=\{p\ge0:\mathtt a^pz_0\sim z_0\}$ とおく．
\begin{enumerate}[label=\textup{(\roman*)},leftmargin=*]
\item ある $p_0\ge1$ について $P_0=\{0\}$ または $P_0=p_0\mathbb Z_{\ge0}$ である．
\item ある $f$ について $l'\mathtt a^fz_0r'\in L$ ならば，
$\{f\ge0:l'\mathtt a^fz_0r'\in L\}$ は，$P_0=\{0\}$ のとき一点であり，そうでないとき $p_0$ を法とする剰余類全体である．
\item $P_0\ne\{0\}$ ならば $\mathtt a\sim\mathtt a^{1+p_0}$ であり，したがってすべての $\mathtt a^k$ は $1\le k'\le p_0$ をみたすある $\mathtt a^{k'}$ と $\equiv_{L,s}$ で同値である．
\end{enumerate}
\end{lemma}

\begin{proof}
語 $\mathtt a^fz_0$ はすべて $\mathtt a$ 混合なので，(I2) から一つの型をもつ．
(i) $P_0$ は和について閉じている．
$P_0$ の $q>p>0$ について，差 $q-p$ についても閉じている：$\D(\mathtt a^pz_0)=\D(z_0)\ne\emptyset$ なので，
$\mathtt a^p(\mathtt a^{q-p}z_0)\sim\mathtt a^pz_0$ に補題~\ref{lem:app-h-cancel} を適用すればよい；$p=0$ の場合は自明である．
(ii) 二つの元 $f<f'$ から $\mathtt a^fz_0\sim\mathtt a^{f'}z_0$，したがって補題~\ref{lem:app-h-cancel} から $f'-f\in P_0$ を得る；
逆に，この集合は $P_0$ の元を足すことについて閉じ，$\mathtt a^{p_0}z_0\sim z_0$ から $p_0$ を引くことについても閉じる．
(iii) $(l',r')\in\D(z_0)$ について $l'\mathtt a^{p_0}z_0r'$ と $l'\mathtt a^{2p_0}z_0r'$ はともに $L$ に入るので，
(I1) から一つの型をもつ純な語 $\mathtt a$ と $\mathtt a^{1+p_0}$ は文脈 $(l'\mathtt a^{p_0-1},z_0r')$ を共有する．
\end{proof}

$D:=\operatorname{lcm}(d,p_0)$ とし，$P_0=\{0\}$ ならば $D:=\infty$ とする．
補題~\ref{lem:app-h-collapse} と補題~\ref{lem:app-h-period}(ii) から，符号を含む上のような指定された積 $u$ すべてについて
\begin{equation}
\label{eq:app-h-member}
 l_\circ ur_\circ\in L\iff\operatorname{val}(u)\equiv V^*\pmod D
\end{equation}
である（$D=\infty$ のときは等号）；価値と $V^*$ は $d$ の倍数である．
同じ語の二つの注釈は $D$ の倍数だけ異なるので，所属は注釈の選び方によらない．

\subsection{数え上げ}

\begin{lemma}
\label{lem:app-h-count}
$(z_1,e_1),\ldots,(z_M,e_M)$（$M\ge6$）は注釈付きの符号で，$|z_j|\le\ell_z$，その価値 $w_j$ の部分和は相異なり，
$\Sigma_w:=\sum_jw_j<D$ とし，$Q$ は $\Sigma_w+\operatorname{val}(Q)=V^*$ をみたす符号と塊の固定した指定された積とする．
$m:=|l_\circ|+|r_\circ|+|Q|+\sum_j|z_j|$ とおく．このとき
\[
 (27/4)^{\lfloor M/3\rfloor}\le(M+1)^2k_V(m+1)^2(M\ell_z+1)^2.
\]
\end{lemma}

\begin{proof}
置換 $\pi$ について，語 $w_\pi:=l_\circ z_{\pi(1)}\cdots z_{\pi(M)}Qr_\circ$ は \eqref{eq:app-h-member} から $L$ に入る．
$G'$ での $w_\pi$ の導出木をとり，単位節点を縮め，各節点に「符号 $z_{\pi(1)},\ldots,z_{\pi(M)}$ のうち最初の文字がその生成語の中にあるものの個数」という重みを与える．
葉の重みは $\le1$，根の重みは $M$ である；重い方の子へ下りると，重み $\kappa\in(M/3,2M/3]$ の節点 $u$ に着く．そのラベルを $X$，生成語の区間を $[p,q)$ とする．
$z_j$ が $[p,q)$ の内部にあるような添字 $j$ の集合を $I(\pi)$ とする；$\kappa-1\le|I(\pi)|\le\kappa$ である．
$\pi$ の\emph{記録}は，$X$，$p$，$q$，および $[p,q)$ の端のうち $M$ 個の符号のどれかの真の内部にあるものそれぞれについての，その符号の添字とずれからなる；
記録は高々 $k_V(m+1)^2(M\ell_z+1)^2$ 通りである．
三つの塊 $l_\circ$，$z_{\pi(1)}\cdots z_{\pi(M)}$，$Qr_\circ$ は，すべての $\pi$ で同じ位置を占める．
最初の文字が $[p,q)$ の中にある符号は，$I(\pi)$ の符号と，内部で始まり $q$ の後で終わる高々一つの符号である；$p$ を通る符号は $p$ の前で始まる．
これが $\kappa-1\le|I(\pi)|$ を与える．$[p,q)$ の各端は，$l_\circ$ の中，$Qr_\circ$ の中，二つの符号の境目，あるいは符号の真の内部にある；
最後の場合にだけそこで符号が切れ，そのときその添字とずれは記録に含まれる．
したがって記録は，各端でこの四つの場合のどれが起こるかも定める．
両端で切れる符号はない．そうなら $[p,q)$ は一つの符号の内部にあり $\kappa\le1$ となるからである．
符号は互いに異なる語である：$z_j=z_k$ ならばその注釈は $D$ の倍数だけ異なるので $w_j\equiv w_k\pmod D$ であり，両者は $[0,D)$ にあるので $w_j=w_k$，部分和が互いに異なるので $j=k$ である．（以下の貼り合わせは位置で行うので，この事実は使わない．）

$\pi,\pi'$ が同じ記録をもつとする．
$w_\pi$ の木の $u$ の部分木を $w_{\pi'}$ の対応する部分木で置き換えると，
$w_\pi[0,p)\,w_{\pi'}[p,q)\,w_\pi[q,|w_\pi|)\in L$ の導出が得られる．
符号の内部にある端は両側で同じ添字とずれをもつので，断片は同じ符号の完全な写しにつながり，
この語は $l_\circ u''Qr_\circ$ となる．ここで $u''$ は完全な符号の積で，$z_j$ はその中に
$1-\mathbf 1[j\in I(\pi)]+\mathbf 1[j\in I(\pi')]$ 回現れる；
これは端が $l_\circ$ や $Qr_\circ$ の中にある場合も含み，$\kappa\ge2$ なので両端が $Qr_\circ$ の中にあることはない．
\eqref{eq:app-h-member} から $\sum_{I(\pi')}w_j\equiv\sum_{I(\pi)}w_j\pmod D$ である；
両辺の和は $[0,\Sigma_w)\subseteq[0,D)$ にあるので等しく，$I(\pi')=I(\pi)$ である．

したがって記録と $I=I(\pi)$（$|I|=j$）は，内側の順序（$\le j!$ 通り）とこの塊の周りの残りの符号の並べ方（$\le(M-j+1)!$ 通り）を除いて $\pi$ を定める．
$j\in\{\kappa-1,\kappa\}$ かつ $M/3<\kappa\le2M/3$ なので $j$ は $J:=[\lfloor M/3\rfloor,\lfloor2M/3\rfloor]$ に入り，
ある $j\in J$ について $\binom{M+1}{j}\le(M+1)k_V(m+1)^2(M\ell_z+1)^2$ である．
$k:=\lfloor M/3\rfloor$ とおく．$\binom{M+1}{\cdot}$ は対称で単峰であり $\lfloor M/3\rfloor+\lfloor2M/3\rfloor\le M$ なので，
すべての $j\in J$ は $\binom{M+1}{j}\ge\binom{M+1}{k}\ge\binom{3k}{k}$ をみたす．
最後に，$0\le i\le3k$ について $t_i:=\binom{3k}{i}3^{-i}(2/3)^{3k-i}$ の和は $1$ であり，
$t_{i+1}/t_i=(3k-i)/(2i+2)$ が $1$ 以上となるのは $i\le k-1$ のとき，かつそのときに限る；
したがって $t_k$ はそれらの最大であり $t_k\ge1/(3k+1)\ge1/(M+1)$ である．
$t_k=\binom{3k}{k}4^k/27^k$ なので，これは $\binom{3k}{k}\ge(27/4)^k/(M+1)$ を与え，補題が従う．
\end{proof}

\subsection{選択と埋め草}

$\Lambda:=\min(V^*,D)$ とおく．したがって $\Lambda\le V^*\le n$ である．

\begin{lemma}[選択]
\label{lem:app-h-select}
値 $w_t:=v_{j_t}$ が超増加（$w_t>w_1+\cdots+w_{t-1}$）であり，$w_1+\cdots+w_M<\Lambda$ かつ $\Lambda\le3^M(v_{i^*}+d)$ をみたす
$M\ge0$ 個の添字 $i^*=j_1>j_2>\cdots>j_M\ge i_0$ がある．
\end{lemma}

\begin{proof}
$v_{i^*}\ge\Lambda$ ならば $M:=0$ ととる．そうでなければ $j_1:=i^*$ ととる．
$j_t$ と $S_t:=w_1+\cdots+w_t<\Lambda$ が与えられたとき，$j_t$ より下で $v_i>S_t$ をみたす最大の添字を $i$ とする；
$v_0=n\ge\Lambda$ なのでそれは存在する（$i_0\ge1$ を思い出す）．
$S_t+v_i\ge\Lambda$ ならば $M:=t$ として止める；そうでなければ $j_{t+1}:=i$ とおく．
値は超増加であり，その和は作り方から $\Lambda$ 未満である．
とられる値は $\Lambda\le n$ 未満であるが，$i<i_0$ については $v_i=n$ である；したがってとられる添字はすべて $i_0$ 以上である．
$i$ の選び方から $v_{i+1}\le S_t$ である：$i+1<j_t$ ならば最大性から，$i+1=j_t$ ならば $v_{j_t}=w_t$ だからである．
\eqref{eq:app-h-star} から $v_i\le2S_t+2d$ である．
したがって $i$ がとられるときは $S_{t+1}+d\le3(S_t+d)$ であり，過程が止まるときは $\Lambda\le S_M+v_i<3(S_M+d)$ である．
$S_1=v_{i^*}$ から $\Lambda<3^M(v_{i^*}+d)$ を得る．
\end{proof}

\begin{lemma}[埋め草]
\label{lem:app-h-filler}
$0<X<n$ を $d$ の倍数とする．
$\hat z_{i_0},\ldots,\hat z_{i^*}$ のうちの高々 $3k_V$ 個の符号と，$d\mid X'$，$X'<v_{i^*}$ をみたす高々一つの純な塊 $\mathtt a^{X'}$ からなる指定された積 $Q$ で，
$\operatorname{val}(Q)=X$ かつ $|Q|\le(3k_V+2)\ell_z$ をみたすものがある．ここで $\ell_z:=5k_V^2\theta$ である．
\end{lemma}

\begin{proof}
$i=i_0,\ldots,i^*$ について順に，$\hat z_i$ を $q_i:=\lfloor X/v_i\rfloor$ 個加えて $X$ を $X-q_iv_i$ で置き換える．
$X<n\le2v_{i_0}$ なので $q_{i_0}\le1$ である．
$i>i_0$ については現在の $X$ は $v_{i-1}$ 未満であり，\eqref{eq:app-h-star} から $v_{i-1}\le2v_i+2d<4v_i$ なので $q_i\le3$ である．
最後の余り $X'<v_{i^*}$ は，すべての価値が $d$ の倍数なので $d$ の倍数であり，塊として加える．
$i^*-i_0\le k_V-3$ なので，符号は高々 $1+3(k_V-3)<3k_V$ 個で，それぞれ長さ $\le\theta+2d\le\ell_z$ であり，
$X'<v_{i^*}\le B_0+2d\le(6k_V^2+4k_V)\theta\le2\ell_z$ である．
\end{proof}

\subsection{定理~\ref{thm:two-letter-fanout-one} の証明}

$A$ が純な生成語をもつとする．(I1) から型 $s(\mathtt a)$ の $A$ の生成語はその純な生成語なので $\tau(A,s(\mathtt a))\le n$ であり，
$\tau(A,s(\mathtt a))\le11\cdot3^{M^*}k_V^2\theta$ を示せば十分である．ここで
\[
 M^*:=\max\bigl(\{5\}\cup\{M\ge6:(27/4)^{\lfloor M/3\rfloor}
 \le(M+1)^4k_V(3k_V+M+4)^2\ell_z^{\,4}\}\bigr)
\]
であり，左辺は $M$ について指数的に増加するのでこれは有限である．
根の部分木の種が高々 $k_V$ 個か，$n_{X_{i^*}}=n$ ならば $n\le2k_V(k_V+2)\theta$ である．
そうでなければ補題~\ref{lem:app-h-select} の添字をとる．

\emph{数え上げ．}
$M\ge6$ とする．$X:=V^*-\sum_tw_t$ とおく；$\sum_tw_t<\Lambda\le V^*\le n$ かつ $w_1>0$ なので $0<X<n$ であり，$X$ は $d$ の倍数である．
$Q$ を補題~\ref{lem:app-h-filler} で与えられるものとする．
符号 $\hat z_{j_1},\ldots,\hat z_{j_M}$ は長さ $\le\ell_z$ で価値が超増加なので，その部分和は相異なり，$\sum_tw_t<\Lambda\le D$ である．
$|l_\circ|+|r_\circ|\le\theta+2d\le\ell_z$ から $m\le(3k_V+M+3)\ell_z$ を得，
$(m+1)^2(M\ell_z+1)^2\le(3k_V+M+4)^2(M+1)^2\ell_z^{\,4}$ とともに補題~\ref{lem:app-h-count} は $M\le M^*$ を与える．
したがってどの場合も $M\le M^*$ であり，補題~\ref{lem:app-h-select} から
\[
 \Lambda\le3^{M^*}(v_{i^*}+d)\le3^{M^*}(B_0+3d)\le10\cdot3^{M^*}k_V^2\theta
\]
である．

\emph{結び．}
$\Lambda=V^*$ ならば $n=V^*+\rho+\rho'<\Lambda+2d$ であり $\tau(A,s(\mathtt a))\le n$ である．
$\Lambda=D<V^*$ ならば $P_0\ne\{0\}$ であり，補題~\ref{lem:app-h-period}(iii) から $n'\le p_0\le D$ について $\mathtt a^n\equiv_{L,s}\mathtt a^{n'}$ となるので，
$\tau(A,s(\mathtt a))\le n'$ である．
小さい場合と合わせて，$\tau(A,s(\mathtt a))\le\max\{2k_V(k_V+2)\theta,\,10\cdot3^{M^*}k_V^2\theta+4k_V^2\theta\}\le11\cdot3^{M^*}k_V^2\theta$ である．

\emph{指数．}
$\mu:=(27/4)^{1/3}$，$\gamma:=\ln3/\ln\mu<1.7260$，$\eta:=1/80$，$C_\eta:=\sup_{M\ge0}(M+1)^6\mu^{-\eta M}<e^{34}$ とおく
（$x>0$ についての $6\ln x-(x-1)\eta\ln\mu$ の最大は $x=6/(\eta\ln\mu)$ でとられる）．
$M^*$ を定める集合の $M\ge6$ について，$\mu^{M-2}\le(27/4)^{\lfloor M/3\rfloor}$ と $3k_V+M+4\le5k_V(M+1)$ から
\[
 \mu^M\le15625\,\mu^2(M+1)^6k_V^{11}\theta^4,\qquad\text{したがって}\qquad
 \mu^{(1-\eta)M}\le15625\,\mu^2C_\eta k_V^{11}\theta^4
\]
を得，$3^M=\mu^{\gamma M}\le(15625\,\mu^2C_\eta k_V^{11}\theta^4)^{\gamma/(1-\eta)}$ である；これは $M=5$ についても成り立つ．
$\gamma/(1-\eta)=80\gamma/79<1.7479$ なので
\[
 \tau(A,s(\mathtt a))\le11\cdot3^{M^*}k_V^2\theta\le C_0k_V^{2+19.23}\theta^{1+6.992}
 \le C_0k_V^{22}\theta^8,\qquad C_0:=11(15625\,\mu^2C_\eta)^{80\gamma/79}<10^{36}
\]
を得る．
$k_V\le2\|G\|$，$\theta\le\bar\theta_G$ であり，$G$ の非終端記号は同じ生成語をもつ $G'$ の非終端記号なので，
$C:=2^{22}C_0<10^{43}$ として定理が従う．$\eta\to0$ とすると，指数は $2+11\gamma<20.99$ と $1+4\gamma<7.91$ にいくらでも近くとれる．\qed

\subsection{他の型と系~\ref{cor:two-letter-fanout-one} の証明}

\begin{lemma}
\label{lem:app-h-other}
$G'$ を上のとおりとし，$A$ を $G'$ の非終端記号とする．
\begin{enumerate}[label=\textup{(\roman*)},leftmargin=*]
\item $s:\Sigma^+\twoheadrightarrow U_1$ で $\mathtt a$ が型 $1$ の唯一の文字であり，$A$ が型 $\zeta$ の生成語をもてば，
$\tau(A,\zeta)\le(k_V-1)\theta+1$ である．
\item $\Sigma=\{\mathtt a,\mathtt b\}$，$s=c$ で，$A$ が内容 $\Sigma$ の生成語をもてば，$\tau(A,\Sigma)\le3k_V\theta+2$ である．
\end{enumerate}
\end{lemma}

\begin{proof}
(i) 型 $\zeta$ の $A$ の生成語の導出で節点の数が最小のものをとり，文字が $\mathtt a$ でない葉に印をつける．
根から印をつけた葉への道の上の節点の部分木を，同じ道の上の同じラベルのより低い節点の部分木で置き換えても，印をつけた葉は残り，したがって型も残る；
よってこの道でラベルは繰り返さず，道は高々 $k_V$ 個の非終端記号の節点をもつ．
そこからぶら下がる高々 $k_V-1$ 個の部分木を長さ $\le\theta$ の生成語で置き換えると，型 $\zeta$ で長さ $\le(k_V-1)\theta+1$ の $A$ の生成語が得られる．
(ii) 内容 $\Sigma$ の生成語の導出で節点の数が最小のものをとり，$\mathtt a$ の葉と $\mathtt b$ の葉に印をつける．
根，分岐点，二つの葉の間の三つの道の区間のそれぞれの中で（異なる区間の間ではなく）同じ縮約を行うと，骨格の節点は高々 $3k_V$ 個である；
そこからぶら下がる高々 $3k_V$ 個の部分木を長さ $\le\theta$ の生成語で置き換えると，長さ $\le3k_V\theta+2$ の内容 $\Sigma$ の生成語が得られる．
どちらの場合も新しい語は同じ型の $A$ の生成語であり，したがって同じ類に入る（補題~\ref{lem:typed-homogeneity}）．
\end{proof}

\begin{proof}[系~\ref{cor:two-letter-fanout-one} の証明]
どちらの場合も $G'$ は $L\cap\Sigma^+$ のファンアウト1，ランク $2\le r_{\max}$ の reduced good な表示で $\bar\theta_{G'}\le\bar\theta_G$ であり，
系~\ref{cor:distribution-representatives} を $G'$ に適用すると
$\|K\|_+\le N_{\mathrm{typ}}(G',s)(1+2\bar\theta_G+r_{\max}\tau(G',s))+\mathbf 1_{\{\eps\in L\}}$ を得る．
残るのは $\tau(G',s)$ と $N_{\mathrm{typ}}(G',s)$ の評価である．

(i) $s$ の正の像は半束の中で $s(\mathtt a)$ と $s(\mathtt b)$ によって生成されるので，$s(w)$ は $c(w)$ のみによる；
したがって $s\preceq c$，$\Cmcf{1}{s}\subseteq\Cmcf{1}{c}$ であり，$\equiv_{L,c}$ は $\equiv_{L,s}$ を細分する．
$G'$ の型付き非終端記号 $A_t$ と，その生成語の一つの $c$ の型 $t'$ について，$\tau_{A_t}(G',s)\le\tau(A,t')$ である．
両方の文字が $c$ について孤立しているので，$\mathtt a$ と $\mathtt b$ のそれぞれについて $c$ に適用した定理~\ref{thm:two-letter-fanout-one} と，
補題~\ref{lem:app-h-other}(ii) から $\tau(G',s)\le C(\|G\|+2)^{22}\bar\theta_G^{\,8}$ である．
$G'$ の型付きの細分は高々 $3k_V$ 個の型付き非終端記号と，規則一つあたり高々9個の型付き規則をもつので，
$N_{\mathrm{typ}}(G',s)\le3k_V+9|P'|\le12(\|G\|+2)$ である．

(ii) $G'$ の型付き非終端記号は，$A$ が $\mathtt a^+$ の生成語をもつ $A_1$ と，$A_\zeta$ の形である．
定理~\ref{thm:two-letter-fanout-one} と補題~\ref{lem:app-h-other}(i) から $\tau(G',s)\le C(\|G\|+2)^{22}\bar\theta_G^{\,8}$ である．
型付き非終端記号は高々 $2k_V$ 個，型付き規則は規則一つあたり高々4個なので
$N_{\mathrm{typ}}(G',s)\le2k_V+4|P'|\le6(\|G\|+\sigma_G)\le12\|G\|$ である．

どちらの場合も $1+2\bar\theta_G\le3\bar\theta_G$ から，$C':=12(C+2)<10^{45}$ として評価を得る．
\end{proof}

\section{ファンアウト1での \texorpdfstring{$U_1$}{U1} の上への型付け}
\label{app:u1-typings}

\subsection{有限の代入可能な言語の文字が互いに素な像}

\begin{lemma}\label{lem:app-u1-insertion}
$L\subseteq\Delta^+$ は有限で Clark--Eyraud 代入可能とする．$PQ\in L$ かつ $PTQ\in L$ となる
$P,Q\in\Delta^*$ と $T\in\Delta^+$ は存在しない．
\end{lemma}

\begin{proof}
そうでないとする．$P\ne\eps$ ならば，語 $P$ と $PT$ は文脈 $(\eps,Q)$ を共有するので
$\D_L(P)=\D_L(PT)$ である；すると $PT^jQ\in L$ から $PT^{j+1}Q\in L$ が従い，
$L$ はあらゆる長さの語 $PT^jQ$ を含み，有限性に反する．$P=\eps$ ならば，
文脈 $(\eps,\eps)$ を共有する $Q$ と $TQ$ について同様に論じる．
\end{proof}

\begin{lemma}\label{lem:app-u1-code}
$h:\Delta^*\to\Sigma^*$ を，すべての $h(\delta)$（$\delta\in\Delta$）が空でなく，
すべての $h(\delta)$ の連接の中で同じ文字が二度現れないような準同型とする．
$L\subseteq\Delta^+$ が有限で Clark--Eyraud 代入可能ならば，$h(L)$ もそうである．
\end{lemma}

\begin{proof}
ある $h(\delta)$ に現れる各文字は，$\delta$ と $h(\delta)$ の中でのその位置を定める．
したがって $h$ は単射であり，語 $h(D)$ の各文字の出現は，符号語 $h(\delta)$ の定まった出現に属する；
隣り合う二つの文字が一つの符号語の出現の中にあるか，ある文字が符号語の出現を始めるか終えるかは，文字そのものから読み取れる．

$x,y\in\Sigma^+$，$x\ne y$ で，$h(L)$ の中で $pxq=h(D_1)$，$pyq=h(D_2)$ とする．
$p$ の内部にある符号語の出現と $q$ の内部にあるものは，両方の語で同じである；
したがって $D_i=PM_iQ$ と書ける．ここで $M_i$ はそれぞれ $x$，$y$ と交わる符号語の出現からなる．
$p$ が符号語の出現 $h(\delta)$ の内部で終われば $M_1$ と $M_2$ はともに $\delta$ で始まり，
$q$ が符号語の出現 $h(\delta')$ の内部で始まれば両者は $\delta'$ で終わる．

$x$ が一つの符号語の出現 $h(\delta)$ の内部にあり，その全体ではないとする；例えば $p$ がその内部で終わるとする（他の場合は対称である）．
このとき $M_1=\delta$ であり，$M_2$ は $\delta$ で始まる；そうでなければ $y=x$ となるので $M_2\ne M_1$ であり，
$T\in\Delta^+$ として $M_2=\delta T$ である．$D_1=P\delta Q$，$D_2=P\delta TQ$ は補題~\ref{lem:app-u1-insertion} に反する．
$x$ と $y$ を入れ替えても同じである．したがって，どちらも符号語全体でない限り一つの符号語の出現の内部にはない．
このとき $M_1=[\delta]X[\delta']$，$M_2=[\delta]Y[\delta']$ である．ここで $X,Y\in\Delta^*$ はそれぞれ $x$，$y$ の内部にある符号語であり，
括弧の文字は対応する切れ目が符号語の内部にあるとき，かつそのときに限り現れる．
$X=\eps\ne Y$ または $Y=\eps\ne X$ ならば，再び補題~\ref{lem:app-u1-insertion} に反し，$X=Y=\eps$ ならば $x=y$ となる．
したがって $X,Y\in\Delta^+$ であり，
\[
 x=\alpha\,h(X)\,\beta,\qquad y=\alpha\,h(Y)\,\beta
\]
である．ここで $\alpha$ は $p$ の後の $h(\delta)$ の部分（または $\eps$），$\beta$ は $q$ の前の $h(\delta')$ の部分（または $\eps$）である．
語 $X$ と $Y$ は $L$ の中で文脈を共有するので $\D_L(X)=\D_L(Y)$ である．

次に $p'xq'\in h(L)$ とする．$\alpha$ と $h(X)$ の間の切れ目，$h(X)$ と $\beta$ の間の切れ目は符号語の出現の境界であり，
それは $x$ の文字から読み取れる；したがってある $P',Q'$ について $p'\alpha=h(P')$，$\beta q'=h(Q')$ であり，$P'XQ'\in L$ である．
すると $P'YQ'\in L$ であり，$h(P'YQ')=p'yq'$ である．対称性から $\D_{h(L)}(x)=\D_{h(L)}(y)$ である．
\end{proof}


\subsection{二文字の型付けについての2次のデータ}

命題~\ref{prop:binary-u1-quadratic} を証明する．
$k:=2^m$，$M:=2^n$，$\psi(w):=|w|_{\mathtt b}+|w|_{\mathtt a}/k$ とおく．$\psi$ は加法的である．
語 $w$ について，$\lambda(w)$ と $\varrho(w)$ を $\mathtt a^*$ に入るその最長の接尾辞と最長の接頭辞の長さとする．

\begin{proof}[命題~\ref{prop:binary-u1-quadratic} の証明]
(i) $\Sigma^+$ の語が $L_{m,n}$ に入るのは，その $\mathtt a$ の極大の連の長さがすべて $k$ で割り切れ，$\psi(w)=M$ であるとき，かつそのときに限る：
そのような語は塊 $\mathtt a^k$ と $\mathtt b$ に一意に分解され，$\psi$ は塊を数える．
$x,y\in\Sigma^+$ が同じ型をもち $uxv,uyv\in L_{m,n}$ とする．加法性から $\psi(x)=\psi(y)$ である．
$x,y$ が純ならば，これは $x=y$ を与える．
そうでなければ，$x'$ を $\mathtt b$ で始まり $\mathtt b$ で終わる語として $x=\mathtt a^{\alpha}x'\mathtt a^{\beta}$ と書き，
同様に $y=\mathtt a^{\alpha'}y'\mathtt a^{\beta'}$ と書く．
$x'$ と $y'$ の内部の $\mathtt a$ の極大の連は $uxv$ と $uyv$ の極大の連なので，その長さは $k$ で割り切れる；
$x$ の両端を通る連の長さは $\lambda(u)+\alpha$ と $\beta+\varrho(v)$ なので，$\alpha\equiv\alpha'$，$\beta\equiv\beta'\pmod k$ である．
任意の文脈 $(p,q)$ について，$pxq$ の $\mathtt a$ の極大の連は，$p$ の末尾にないもの，長さ $\lambda(p)+\alpha$ の連，
$x'$ の内部の連，長さ $\beta+\varrho(q)$ の連，$q$ の先頭にないもの，である；
$pyq$ についても $\alpha',\beta',y'$ で同じことが成り立つ．
したがって割り切れる条件は一致し，$\psi(pxq)=\psi(pyq)$ であり，$pxq\in L_{m,n}\iff pyq\in L_{m,n}$ である．
有限言語は文脈自由である．

(ii) 非終端記号 $C_0,\ldots,C_m$，$A_0,\ldots,A_n$，開始記号 $A_n$ と，
規則 $C_0\to\mathtt a$，$C_j\to C_{j-1}C_{j-1}$（$1\le j\le m$），
$A_0\to C_m\mid\mathtt b$，$A_i\to A_{i-1}A_{i-1}$（$1\le i\le n$）をとる．
このとき $L(C_j)=\{\mathtt a^{2^j}\}$，$L(A_0)=\{\mathtt a^k,\mathtt b\}$ であり，$L(A_i)$ は
$L(A_0)$ の語 $2^i$ 個の連接の集合なので $L(A_n)=L_{m,n}$ である；文法は good かつ reduced である．
非終端記号は $m+n+2$ 個で，規則の大きさは $2$，$5$（$m$ 個），$3$，$2$，$5$（$n$ 個）である（定義~\ref{def:standard-mcfg}）．
これは $6m+6n+9$ を与える．
ファンアウト1では，規則 $R:B\to\rho(B_1,\ldots,B_q)$ を使う最短の文の長さは
$\theta_G(R)=o_B+c_R+\sum_jd_{B_j}$ である．ここで $d_C$ は $C$ の最短の生成語の長さ，
$o_C$ は文の中での $C$ の出現の周りの文字の最小の個数，$c_R$ は $R$ が書く終端記号の個数である；
ここでは $d_{A_0}=1$，$d_{A_i}=2^i$，$d_{C_j}=2^j$，$o_{A_i}=M-2^i$，$o_{A_0}=o_{C_m}=M-1$，
$o_{C_j}=M-1+k-2^j$ である．
これは $A_i$ の規則と $A_0\to\mathtt b$ について $\theta_G(R)=M$ を，$A_0\to C_m$ と $C_j$ の規則について
$\theta_G(R)=M+k-1$ を与える；したがって $\bar\theta_{G_{m,n}}=2^m+2^n-1$ である．

(iii) $L_{m,n}$ の唯一の純な語は $\mathtt a^{kM}$ なので，命題~\ref{prop:rare-type} からすべての特徴標本はそれを含み，
$\|F\|_+\ge kM+1>2^{m+n}$ である．$m=n$ のとき，$\bar\theta_{G_{n,n}}<2^{n+1}$ から
$\bar\theta_{G_{n,n}}^{\,2}/4<4^n<\|F\|_+$ である．評価 $p(\|G\|)\bar\theta_G$ はすべての $n$ について
$4^n\le p(12n+9)\,2^{n+1}$ を与えるが，これはどの多項式 $p$ についても成り立たない．
\end{proof}

\section{ファンアウト2以上での二文字の内容型付け}
\label{app:two-letter-content}

この付録は定理~\ref{thm:two-letter-content} を証明する．
以下では $\Sigma=\{\mathtt a,\mathtt b\}$，$c:=c_\Sigma$，$f\ge2$ を固定し，$L\in\Cmcf{f}{c}$ とし，
$G=(V,\Sigma,P,S,\fan)$ を $L(G)=L\cap\Sigma^+$ をみたすファンアウト $\le f$，ランク $\le r:=r_{\max}$ の reduced good な文法とする；
$v:=|V|$，$\theta:=\bar\theta_G$，$D_f:=(6f+1)3^f$ とおく．
考えている型付き非終端記号のアリティの条件を使う補題~\ref{lem:app-two-basic}(i) を除き，
議論は $(1,c)$- と $(2,c)$-タプル代入可能性のみを使う．これらは $f\ge2$ なので成り立つ．
各 $A\in V$ について文脈 $E_A$ は補題~\ref{lem:thickness-equivalence}(iv) のものとする．したがって $|E_A|_{\mathrm{fix}}\le\theta$
かつ $E_A[L_A(G)]\subseteq L(G)$ である；各 $B\in V$ について $\|\omega_B\|_1\le\theta$ をみたす $\omega_B\in L_B(G)$ を固定し（項目~(ii)），
すべての規則が高々 $\theta$ 個の終端記号を書くことを思い出す（項目~(i)）．
導出木の中で，節点 $u$ の\emph{生成語}はそれが導くタプル，$\ell_u$ はその総長，$B_u$ はその非終端記号であり，
$\ell_u>\theta$ のとき $u$ は\emph{重い}という；節点はその雛形の終端記号を\emph{書く}．
文字 $x,y$ は常に $\Sigma$ の二つの文字をある順序で表す．
語の\emph{連}は極大の因子 $x^k$（$k\ge1$）である；最初の連でも最後の連でもない連を\emph{内部の}連という．
したがって内部の $y$-連の前後には $x$ がある．
\emph{$x$-骨格} $\sigma_x(w)$ は $w$ の $x$-連をそれぞれ一つの $x$ で置き換えたものであり，
$|w|_x$ が $w$ の $x$-連の個数を超えるとき，$w$ は\emph{余分な} $x$ をもつという．
$L$ のある語が因子 $x^k$ をもつとき \emph{$x^k$ が現れる}といい，
同じ Parikh ベクトル $\Psi(u)=\Psi(w)$ をもつ $u,w\in\Sigma^+$ について $u\in L$ と $w\in L$ が同値であるとき，
$L$ は \emph{Parikh 飽和}であるという．

\paragraph{証明の地図．}
表は，各段階について，それが使うものと，それが使われる場所を示す．
\begin{center}
\footnotesize
\setlength{\tabcolsep}{3pt}
\renewcommand{\arraystretch}{1.15}
\begin{tabularx}{\linewidth}{@{}>{\raggedright\arraybackslash}p{0.24\linewidth}>{\raggedright\arraybackslash}X>{\raggedright\arraybackslash}p{0.26\linewidth}@{}}
\hline
段階 & 使うもの & 使われる場所\\
\hline
補題~\ref{lem:app-two-basic}(i) 型付きの閉包
 & 補題~\ref{lem:typed-homogeneity}，アリティ $d$ の代入可能性
 & 補題~\ref{lem:app-two-saturated}\\
補題~\ref{lem:app-two-basic}(ii) 恒等の道の切断
 & $G$ が good であること
 & (P2)，補題~\ref{lem:app-two-saturated} の混合の場合\\
補題~\ref{lem:app-two-runs} 連の交換
 & $(1,c)$，$(2,c)$-代入可能性，$x^4$ が現れる
 & 補題~\ref{lem:app-two-alternate}，命題~\ref{prop:app-two-nonsaturated}\\
補題~\ref{lem:app-two-alternate}
 & 補題~\ref{lem:app-two-runs}(i),(iii)
 & 命題~\ref{prop:app-two-nonsaturated}\\
補題~\ref{lem:app-two-insertion} 置換が作る連
 & 補題~\ref{lem:witness-length-template}(i)，節点数最小の最短の $T$
 & 命題~\ref{prop:app-two-nonsaturated}\\
命題~\ref{prop:app-two-nonsaturated}（飽和でない）
 & 上の三つの補題；道の上の型付きのラベル
 & 定理~\ref{thm:two-letter-content}\\
補題~\ref{lem:app-two-saturated}，和集合が $\Sigma$
 & 補題~\ref{lem:app-two-basic}；道に沿った床関数
 & 定理~\ref{thm:two-letter-content}\\
(P1) 行の剛性
 & $(1,c)$-代入可能性，$K$ の最小性
 & (P3)，(P4)\\
(P2) 高さ，混合の節点
 & 補題~\ref{lem:app-two-basic}(ii)，$K$ の最小性
 & (P3)，(P6)\\
(P3) 重い純な節点
 & (P1)，(P2)
 & (P4)，(P5)\\
(P4) 共通の比 $\alpha$
 & (P3)
 & (P5)，(P6)\\
(P5) $\alpha$ の評価
 & (P3)，(P4)
 & (P6)\\
(P6) 数え上げ
 & (P2)，(P4)，(P5)
 & 定理~\ref{thm:two-letter-content}\\
\hline
\end{tabularx}
\end{center}

\begin{lemma}\label{lem:app-two-basic}
\textup{(i)} アリティ $d$ のすべての型付き非終端記号 $X=A_{\mathbf t}$ について，
$\operatorname{Cl}_{L,c}(X)=\{\tuple z\in(\Sigma^+)^d:c^{(d)}(\tuple z)=\mathbf t,\
E_A[\tuple z]\in L\}$ である；したがって $\tau_X$ はそのようなタプルの最小の総長である．

\textup{(ii)} $G$ の導出木の中で，$u$ が同じ非終端記号をもつ $u'$ の真の祖先であり，
それらの生成語の総 Parikh ベクトルが等しいとする．
このとき $u$ から $u'$ への道の上のすべての規則はランク1で，終端記号を書かず，成分ごとの恒等写像である．
したがって $u$ の部分木を $u'$ の部分木で置き換えても導かれるタプルは変わらず，少なくとも一つの節点が除かれる．
\end{lemma}

\begin{proof}
(i) $\tuple u\in L_X$ を固定する．補題~\ref{lem:typed-homogeneity} から，
$\operatorname{Cl}_{L,c}(X)$ は $\tuple u$ の型と分布をもつタプルの集合であり，それらは $E_A\in\D_L(\tuple u)$ に受理される．
$E_A$ に受理される型 $\mathbf t$ のタプル $\tuple z$ はこの文脈を $\tuple u$ と共有するので，代入可能性から $\D_L(\tuple z)=\D_L(\tuple u)$ である．

(ii) 線形性と非削除性から，$u'$ の生成語のすべての文字は $u$ の生成語に現れるので，
Parikh ベクトルが等しいことは，道の上で書かれる終端記号と，道から外れる兄弟を排除する．兄弟は空でない生成語をもつからである（$\eps$ なし）．
したがって道の上のすべての規則は，あるファンアウト $d'$ の子を一つもち，その $d'$ 個の変数はそれぞれ一度（非削除性），
順に（非置換性）現れ，親の各成分は空でない（$\eps$ なし）．
一つの成分の中の二つの変数は連続することになり，それは非併合性が禁じる；
したがってすべての成分は一つの変数であり，ファンアウトは一致し，規則は恒等写像である．
\end{proof}

\begin{lemma}[連の交換]\label{lem:app-two-runs}
$x^4$ が現れるとする．
\begin{enumerate}[label=\textup{(\roman*)},leftmargin=*]
\item $w,w'\in\Sigma^+$ が $\sigma_x(w)=\sigma_x(w')$ かつ $|w|_x=|w'|_x$ をみたせば，$w\in L$ と $w'\in L$ は同値である．
\item $L$ のある語が因子 $xxyx$ または $xyxx$ をもてば，$L$ は Parikh 飽和である．
\item $L$ が Parikh 飽和でなければ，因子 $xx$ をもつ $L$ のすべての語の中で，すべての内部の $y$-連の長さは少なくとも $2$ である．
\item $L$ が Parikh 飽和でなく $y^4$ が現れなければ，因子 $xx$ をもつ $L$ のすべての語は高々一つの内部の $y$-連をもつ．
\end{enumerate}
\end{lemma}

\begin{proof}
(i) $c(x)$ は冪等なので，補題~\ref{lem:power-transfer} が $u=x$，$i_0=p=1$，$i=\delta=j=1$ で適用できる：
$x^4$ が現れるので，$w_{\mathrm M}\ne\eps$ ならば
$w_{\mathrm L}xw_{\mathrm M}xxw_{\mathrm R}\in L$ と $w_{\mathrm L}xxw_{\mathrm M}xw_{\mathrm R}\in L$ は同値である．
$R$ と $R'$ を $|R|\ge2$ である語 $w$ の $x$-連とし，例えば $R'$ が $R$ の左にあるとする（他の場合は対称である）．
最初に示した $x$ を $R'$ の最後の文字，$xx$ を $R$ の最初の二文字として $w=U\,x\,Z\,xx\,V$ と書く；
因子 $Z$ は二つの連を隔てるので $Z\ne\eps$ である．
したがって $U\,xx\,Z\,x\,V\in L$ と $w\in L$ は同値である．
これは一つの $x$ を $R$ から $R'$ へ移し，$x$-骨格と $|w|_x$ を保つ．
同じ $x$-骨格をもつ語は同じ個数 $\rho$ の $x$-連をもち，その連の長さは $|w|_x$ の $\rho$ 個の正の部分への合成である；
そのような任意の二つの合成は，大きさ少なくとも $2$ の部分から別の部分へ一単位を移す操作でつながる．

(ii) 例えば $UxxyxV\in L$ とする（他の場合は対称である）．(i) から $UxyxxV\in L$ である．
文脈 $(Ux,xV)$ はどちらも内容 $\Sigma$ の $xy$ と $yx$ を受理するので $\D_L(xy)=\D_L(yx)$ である：
すべての文脈で $xy$ を $yx$ に置き換えることができ，その逆もできる．
隣接する互換は一つの Parikh ベクトルをもつすべての語をつなぐ．

(iii) $w\in L$ が因子 $xx$ と長さ $1$ の内部の $y$-連をもつとする．
このとき $w$ は余分な $x$ をもち，(i) から余分な文字をその $y$-連の前の $x$-連へ移すことができる；
得られる $L$ の語は因子 $xxyx$ をもつので，(ii) から $L$ は Parikh 飽和である．

(iv) $w\in L$ が因子 $xx$ と二つの内部の $y$-連をもつとする．
このとき $w$ は，$x$-連 $R$ で隔てられた連続する二つの内部の $y$-連 $y^e$ と $y^{e'}$ をもつ；$A_0$ を $y^e$ の前の $x$-連とする．
(iii) から $e,e'\ge2$ である．
$w$ は余分な $x$ をもつので，(i) は $w$ の $x$-骨格と文字 $x$ の個数をもつ $L$ の二つの語を与える：
$R$ の長さが $2$ である語 $W_2=U\,x\,y^e\,xx\,V'$（$y^e$ の前に示した $x$ は $A_0$ の最後の文字であり，$V'$ は $y^{e'}$ で始まる）と，
$R$ の長さが $1$ である語 $W_1=U_1\,x\,y^e\,x\,y^{e'}V_1$ である．
文脈 $(U,xV')$ は $W_2$ を与える $xy^ex$ と，$U\,xx\,y^e\,x\,V'$ を与える $xxy^e$ を受理する．
後者は $W_2$ の $x$-骨格と $x$ の個数をもち，(i) から $L$ に入る．
どちらのタプルも内容 $\Sigma$ をもつので $\D_L(xy^ex)=\D_L(xxy^e)$ である．
文脈 $(U_1,y^{e'}V_1)$ は $W_1$ を与える $xy^ex$ を受理するので $xxy^e$ も受理し，$U_1\,xx\,y^{e+e'}V_1\in L$ である．
$e+e'\ge4$ なので，この語は因子 $y^4$ をもち，矛盾である．
\end{proof}

\begin{lemma}\label{lem:app-two-alternate}
$L$ が Parikh 飽和でなく，$\mathtt a^4$ と $\mathtt b^4$ がともに現れるならば，
少なくとも四つの連をもつ $L$ のすべての語のすべての連の長さは $1$ である．
\end{lemma}

\begin{proof}
$w\in L$ が少なくとも四つの連をもつとする．その二番目と三番目の連は内部の連で異なる文字からなるので，
各文字 $y$ について，$w$ は内部の $y$-連 $\beta$ と少なくとも二つの $y$-連をもつ．
$y$ についての補題~\ref{lem:app-two-runs}(i) により，文字 $y$ を $\beta$ から別の $y$-連へ移すと，
$\beta$ の長さが $1$ で $x$-連が $w$ のものである $L$ の語が得られる．
$w$ の $x$-連の長さが少なくとも $2$ ならば，この語は因子 $xx$ と長さ $1$ の内部の $y$-連をもち，
補題~\ref{lem:app-two-runs}(iii) に反する．
したがって $x$ のどちらの選び方についても，$w$ のすべての $x$-連の長さは $1$ である．
\end{proof}

\begin{lemma}[短い置換が作る連]\label{lem:app-two-insertion}
$X=A_{\mathbf t}$ をアリティ $d\le f$ の型付き非終端記号とし，$T$ を，$X$ からの $\widetilde G_0$ の導出で，
その生成語 $\tuple u$ の総長が最小で，その中で節点の数が最小のものとする．
$T$ が $|Q|\ge6f+2$ をみたす重い節点の反鎖 $Q$ をもてば，ある文字 $x$ について次が成り立つ：
\textup{(a)} $\tuple u$ のある成分は内容 $\{x\}$ と少なくとも $5$ の長さをもち，したがって $x^5$ が現れる；
\textup{(b)} $L$ のある語 $W$ は因子 $xx$ と少なくとも四つの連をもち，$\mathtt a^5$ と $\mathtt b^5$ がともに現れるのでなければ，
$W$ は二つの内部の $y$-連をもつ．
\end{lemma}

\begin{proof}
内容 $\{x\}$ の各成分 $u_i$（\emph{純な}成分）の最初と最後の $\min\{2,|u_i|\}$ 文字に印をつけ，
内容 $\Sigma$ の各成分の一つの $\mathtt a$ と一つの $\mathtt b$ に印をつける；印は高々 $4f$ 個である．
各印は一つの節点が書き，$Q$ の節点は互いに素な部分木をもつので，
部分木が印を書かない $Q$ の節点の集合 $Q_0$ は $|Q_0|\ge|Q|-4f\ge2f+2$ をみたす．

$h\in Q_0$ の非終端記号を $C_h$ とする．型の添字を消し，$h$ の部分木を $\omega_{C_h}$ の導出で置き換える．
新しい根のタプルは $L_A(G)$ に入り，より短い．非削除性から $h$ の生成語のすべての文字は $\tuple u$ に現れ，
$\|\omega_{C_h}\|_1\le\theta<\ell_h$ だからである．
印は内容 $\Sigma$ のすべての成分の内容を保つ．
すべての純な成分が純なままならば，新しいタプルは型 $\mathbf t$ をもち，
$\lambda:=\widehat c$ とした補題~\ref{lem:witness-length-template}(i) から $L_X$ に入るが，これは $T$ の選び方に反する．
したがって $\omega_{C_h}$ は，純な成分の中にある $h$ の塊（$h$ の生成語の成分が占める根の成分の区間）に他方の文字を入れる．
塊は印を含まないので，その純な成分の長さは少なくとも $5$ であり，塊はその最初の二文字の後，最後の二文字の前にある．
$E_A[\tuple u]\in L$ なので，これは (a) を示す．

各 $h\in Q_0$ は高々 $f$ 個の塊をもつ．
一つの根の成分の中でそれらの塊が間に文字なしに隣り合うとき，二つの所有者を結ぶ．
各塊は高々二つの隣り合う塊をもつので，このグラフの最大次数は高々 $2f$ であり，
$|I|\ge|Q_0|/(2f+1)>1$，したがって $|I|\ge2$ をみたす独立集合 $I$ をもつ．
すべての $h\in I$ の部分木を同時に置き換えると，$\tuple u'\in L_A(G)$ と $W:=E_A[\tuple u']\in L$ を得る．
文字 $x$ の純な成分を固定し，他方の文字を $y$ と書く．
$I$ の塊の外にあるその文字は変わらず，したがって $x$ に等しい；
これは $I$ の外の所有者の塊（保たれる）と，最初の二文字と最後の二文字（印がついている）を含む．
したがって新しい成分の各極大の $y$-連の前後には $x$ があり，それは間に文字なしに連続する $I$ の塊の和の中にある．
そのような塊は一つの所有者に属する：
間に $I$ の塊がない $I$ の相異なる所有者 $h\ne h'$ の二つの塊は，変わらない文字で隔てられる．
その隙間は空でない（そうでなければ $h$ と $h'$ は結ばれる）か，$I$ の外の所有者の塊を含むからである．
したがって新しい成分の各 $y$-連は $I$ の中に一意な\emph{作る}所有者をもち，前の段落から，各 $h\in I$ は，
その文字を $x_h$ と呼ぶ純な成分の中に少なくとも一つの連を作る．
$I$ の中の $h\ne h'$ について，それらが作る連は異なる：
同じ純な成分の中にあれば作る所有者が異なり，異なる純な成分の中にあれば $\tuple u'$ の異なる成分の中にある．

$I$ の中の $h_1\ne h_2$ をとる．$x_{h_1}=x_{h_2}=:x$ ならば，$W$ は因子 $xx$（長さ少なくとも $5$ の純な成分の最初の二文字）と
二つの内部の $y$-連をもち，したがって少なくとも五つの連をもつ．
そうでなければ，$h_1$ と $h_2$ に (a) の議論を適用して $\mathtt a^5$ と $\mathtt b^5$ はともに現れ，
$x:=x_{h_1}$，$y:=x_{h_2}$ とすると，二つの純な成分はそれぞれ連 $x,y,x$ と $y,x,y$ を含む；
これらは $W$ の六つの相異なる連であり，$W$ は因子 $xx$ をもつ．
\end{proof}

\begin{proposition}\label{prop:app-two-nonsaturated}
$L$ が Parikh 飽和でなければ $t_{\mathrm{typ}}(G,c)\le D_f(r+1)v\theta$ である．
\end{proposition}

\begin{proof}
$X$ と $T$ を補題~\ref{lem:app-two-insertion} のとおりとする．
まず，$T$ の重い節点のすべての反鎖が高々 $6f+1$ 個の節点をもつことを示す．
そうでなければ補題~\ref{lem:app-two-insertion} が適用できる．
$\mathtt a^4$ も $\mathtt b^4$ も現れなければ，(a) は不可能である．
両方が現れれば，(b) の語 $W$ は少なくとも四つの連と長さ少なくとも $2$ の連をもち，補題~\ref{lem:app-two-alternate} に反する．
ちょうど一方，例えば $x_0^4$ が現れれば，(a) から $x=x_0$ であり $y^5$ は現れないので，
(b) から $W$ は因子 $xx$ と二つの内部の $y$-連をもち，補題~\ref{lem:app-two-runs}(iv) に反する．

型付きのラベルは非終端記号とその成分の内容の組なので，高々 $3^fv$ 個ある．
$T$ の道の上で型付きのラベルは繰り返さない．二つの等しいラベルの間で切ると，
生成語が長くならず節点の少ない $X$ からの導出が得られるからである；したがって道は高々 $3^fv$ 個の節点をもつ．
重い節点は祖先について閉じており，一つの深さにあるものは反鎖をなすので，それらは高々 $(6f+1)3^fv=D_fv$ 個である．
根が重くなければ $\|\tuple u\|_1\le\theta$ である．
そうでなければ，$\tuple u$ のすべての文字は重い節点が書くか，その重くない子の生成語の中にあり，
各重い節点は高々 $\theta$ 文字を書き，高々 $r$ 個の子をもつ．したがって $\|\tuple u\|_1\le D_f(r+1)v\theta$ である．
\end{proof}

\begin{lemma}\label{lem:app-two-saturated}
$L$ が Parikh 飽和ならば $\tau(G,c)\le C_{f,r}v^2\theta^4H_\theta$ である．ここで $C_{f,r}:=(r+1)(4r(3f+1)+D_f)+3$ である．
\end{lemma}

\begin{proof}
アリティ $d$ の $X=A_{\mathbf t}$ を固定し，$(p,q):=\Psi(E_A)$ とおく．したがって $p+q\le\theta$ である．
また $\mathcal P:=\Psi(L\cap\Sigma^+)$ とおく．
補題~\ref{lem:app-two-basic}(i) と Parikh 飽和から，$\tau_X$ は
$(p,q)+\sum_i\Psi(z_i)\in\mathcal P$ をみたす型 $\mathbf t$ のタプル $\tuple z$ についての $\|\tuple z\|_1$ の最小値である．
$(x_B,y_B):=\Psi(\omega_B)$ と書き，導出木の節点 $u$ について，その生成語の Parikh ベクトルを $(a_u,b_u)$ とする．
$u$ が語 $w$ の成功した導出の中にあれば，その部分木を $\omega_{B_u}$ の導出で置き換えると，
Parikh ベクトル $\Psi(w)-(a_u,b_u)+(x_{B_u},y_{B_u})$ をもつ $L$ の語が得られる；
$B_{u'}=B_u$ である真の子孫 $u'$ について，$u$ の部分木を $u'$ の部分木で置き換えると，
$\Psi(w)-(a_u-a_{u'},b_u-b_{u'})$ をもつ語が得られる．
以下の比較はすべてこの種類のものである：それは，固定した導出からそのような置換で得られる完全な成功した導出が生成する $L$ の語の
Parikh ベクトルに関するものであり，Parikh 飽和から，語の $L$ への所属はその Parikh ベクトルのみによる．
内容 $t_1,\dots,t_d$ の和集合は $\Sigma$ であるか，そうでなければそれらはすべて一つの単集合に等しい．

\emph{和集合が $\Sigma$ の場合．}
$t_{\mathtt a}$ と $t_{\mathtt b}$ を，内容がそれぞれ $\mathtt a$，$\mathtt b$ を含む成分の個数とする；どちらも $\{1,\dots,f\}$ に入る．
$n_{\mathtt a}$ 個の文字 $\mathtt a$ と $n_{\mathtt b}$ 個の文字 $\mathtt b$ をもつ型 $\mathbf t$ のタプルが存在するのは
$n_{\mathtt a}\ge t_{\mathtt a}$ かつ $n_{\mathtt b}\ge t_{\mathtt b}$ のとき，かつそのときに限るので，
$\tau_X$ はこれらの下界と $(p+n_{\mathtt a},q+n_{\mathtt b})\in\mathcal P$ をみたす最小の $n_{\mathtt a}+n_{\mathtt b}$ である．
最小にする組を固定し，$c_{\mathtt a}:=n_{\mathtt a}-t_{\mathtt a}+1\ge1$，
$c_{\mathtt b}:=n_{\mathtt b}-t_{\mathtt b}+1\ge1$ とおき，
$W:=\mathtt a^{p+n_{\mathtt a}}\mathtt b^{q+n_{\mathtt b}}$ の成功した導出 $T$ で節点の数が最小のものを固定する．
重い節点 $u$ が $a_u<c_{\mathtt a}$ かつ $b_u<c_{\mathtt b}$ をみたせば，その置換は
$n'_{\mathtt a}\ge t_{\mathtt a}$，$n'_{\mathtt b}\ge t_{\mathtt b}$，
$n'_{\mathtt a}+n'_{\mathtt b}<n_{\mathtt a}+n_{\mathtt b}$ をみたす Parikh ベクトル $(p+n'_{\mathtt a},q+n'_{\mathtt b})\in\mathcal P$ を与える；
したがってすべての重い節点は $a_u\ge c_{\mathtt a}$ または $b_u\ge c_{\mathtt b}$ をみたす．
$u'$ が同じ非終端記号をもつ $u$ の真の子孫ならば，$a_u-a_{u'}\ge c_{\mathtt a}$ または $b_u-b_{u'}\ge c_{\mathtt b}$ である：
差が零ならば節点数の最小性に反し（補題~\ref{lem:app-two-basic}(ii)），他の小さい差は組の選び方に反する．
道に沿って $(\lfloor a_u/c_{\mathtt a}\rfloor,\lfloor b_u/c_{\mathtt b}\rfloor)$ は成分ごとに非増加であり，
それが一定の間は非終端記号は繰り返さない．
$n\ge t\ge1$，$p\ge0$ について $\lfloor(p+n)/(n-t+1)\rfloor\le p+t$ である
（$(p+t)(n-t+1)-(p+n)=(n-t)(p+t-1)\ge0$）ので，それは高々 $1+(p+t_{\mathtt a})+(q+t_{\mathtt b})\le\theta+2f+1$ 個の値をとり，
道は高々 $v(\theta+2f+1)$ 個の節点をもつ．
一つの深さにある重い節点は互いに素な生成語をもつので，それらは高々
$(p+n_{\mathtt a})/c_{\mathtt a}+(q+n_{\mathtt b})/c_{\mathtt b}\le\theta+2f$ 個である．
重い節点は祖先について閉じており，すべての文字は重い節点が書くか，その高々 $r$ 個の重くない子の一つの中にある；したがって
$\tau_X\le|W|\le\max\{\theta,(r+1)\theta(\theta+2f)(\theta+2f+1)v\}\le(2f+1)(2f+2)(r+1)v\theta^3$ である．

\emph{すべての内容が $\{\mathtt a\}$ に等しい場合}（$\{\mathtt b\}$ の場合は対称である）．$h:=d\in\{1,\dots,f\}$ とする．
$K:=\min\{k\ge h:(p+k,q)\in\mathcal P\}$ とおく；このとき $\tau_X=K$ であり，
そうでなければ $\tau_X\le(2f-1)\theta$ なので，$K>\theta+2(f-1)\ge p+2(h-1)$ と仮定してよい．
$W:=\mathtt b^q\mathtt a^{p+K}$ の成功した導出 $T$ で節点の数が最小のものを固定する．

(P1) \emph{行の剛性}：$c_0\ge1$，$1\le\delta\le K-h$ をみたす $(c_0,y),(c_0+\delta,y)\in\mathcal P$ は存在しない．
そうでなければ，文脈 $(\mathtt b^y\mathtt a^{c_0-1},\eps)$ は $\mathtt a$ と $\mathtt a^{1+\delta}$ を受理するので，
$(1,c)$-代入可能性から $\D_L(\mathtt a)=\D_L(\mathtt a^{1+\delta})$ である（どちらも内容 $\{\mathtt a\}$ をもつ）；
文脈 $(\mathtt b^q\mathtt a^{p+K-\delta-1},\eps)$ は $\mathtt a^{1+\delta}$ を受理するので $\mathtt a$ も受理し，
$h\le K-\delta<K$ について $(p+K-\delta,q)\in\mathcal P$ となり，$K$ の選び方に反する．

(P2) \emph{高さ．}
$u'$ が同じ非終端記号をもつ $u$ の真の子孫で $b_u=b_{u'}$ ならば $a_u-a_{u'}\ge K-h+1$ である：
値 $0$ は節点数の最小性に反し（補題~\ref{lem:app-two-basic}(ii)），$[1,K-h]$ の値は $K$ の選び方に反する．
道を下ると $b_u$ は非増加である（非削除性から子の生成語は親の生成語の中に現れる）ので，
$b_u$ の一つの値をもつ道の節点は区間をなす；$b_u$ は高々 $q+1$ 個の値をとり，
一つの非終端記号と $b_u$ の一つの値をもつ三つの節点は $p+K\ge a_u\ge2(K-h+1)$ を与えるが，$K>p+2h-2$ なのでこれは偽である．
したがって道は高々 $2v(q+1)$ 個の節点をもつ．
生成語が $\mathtt b$ を含む節点を\emph{混合}という．
$1\le k\le q$ について，$b_u=k$ をみたし，そのような真の祖先をもたない混合の節点は互いに素な生成語をもつので，
高々 $\lfloor q/k\rfloor$ 個である；そのうちの一つの下では，そのような二つの節点は比較可能であり（そうでなければその生成語は $2k$ 個の文字 $\mathtt b$ を含む），
したがって鎖をなし，それは上の議論から高々 $2v$ 個の節点をもつ．したがって混合の節点は
$M\le2v\sum_{k=1}^q\lfloor q/k\rfloor\le2vqH_q\le2v\theta H_\theta$ 個である．

(P3) \emph{重い純な節点．}
混合でない節点 $u$ で $\ell_u=a_u>\theta+f-1$ をみたすものを\emph{重い純な}節点と呼び，
$g_u:=\ell_u-x_{B_u}\ge f$ とおく．
比較可能な二つの重い純な節点は異なる非終端記号をもつ．そうでなければ (P2) から上の節点は
$\ell_u\ge K-h+1+\theta+f>p+K=|W|_{\mathtt a}$ をもつからである．
特に，一つの非終端記号をもつ重い純な節点は互いに素な生成語をもつ．
重い純な節点の反鎖 $J$ について，そのすべてを置き換えると Parikh ベクトル $(p+K-\sigma(J),q+Y(J))$ をもつ語が得られる．
ここで $\sigma(J):=\sum_{u\in J}g_u$，$Y(J):=\sum_{u\in J}y_{B_u}$ である；
この語は完全な成功した導出が生成するので，その Parikh ベクトルは $\mathcal P$ に入る．
したがって $\sigma(J),\sigma(J')\le K-h$ かつ $Y(J)=Y(J')=:Y$ をみたす二つの反鎖 $J,J'$ は
$(p+K-\sigma(J),q+Y),(p+K-\sigma(J'),q+Y)\in\mathcal P$ を与える．
その第一座標は少なくとも $p+h\ge1$ で，差は $|\sigma(J)-\sigma(J')|\le K-h$ である；(P1) から $\sigma(J)=\sigma(J')$ である．

(P4) \emph{繰り返す非終端記号．}
$B$ が $m_B\ge2$ 個の重い純な節点のラベルならば，(P3) からそれらの生成語は互いに素であり，
これらの節点は互いに比較不能で，そのどの集合も反鎖である．したがって各々は $\ell_u\le p+K-(\theta+f)\le K-f$，
$g_u\le K-f\le K-h$ をみたす；(P3) で二つの単集合を比べると共通の値 $g_B$ を得，したがって共通の $\ell_B$ を得る．
さらに $y_B=0$ ならば，一つの置換は $h\le K-g_B<K$ について $(p+K-g_B,q)\in\mathcal P$ を与え，これは不可能である；したがって $y_B\ge1$ である．
$m_B>\theta$ のとき $B$ を\emph{頻出}という．
頻出の $B,C$ について，$B$ をラベルとする $y_C$ 個の節点と，別に $C$ をラベルとする $y_B$ 個の節点を選ぶ．
$m_B>\theta\ge y_C$，$m_C>\theta\ge y_B$ なのでこれは可能である；
どちらの集合も $Y=y_By_C$ をもつ反鎖である．$y_C\le m_B-1$，$g_B\le\ell_B$ であり，$m_B$ 個の互いに素な生成語から $m_B\ell_B\le p+K$ なので，$y_Cg_B\le(m_B-1)\ell_B\le p+K-\ell_B\le K-f\le K-h$ である（$\ell_B\ge\theta+f\ge p+f$ による）．$C$ についても同様である．
(P3) から $y_Cg_B=y_Bg_C$ なので，すべての頻出の $B$ は共通の比 $\alpha:=g_B/y_B$ をもち，$\ell_B=x_B+\alpha y_B\le(1+\alpha)\theta$ である．

(P5) \emph{比．}
$\alpha_*:=(r+1)(\theta+f-1)(2\theta+1)v\le3f(r+1)v\theta^2$ とおく；
頻出の非終端記号が存在するとき $\alpha\le\alpha_*$ を示す．$\alpha\le\theta$ ならば明らかである．
そうでなければ，頻出の非終端記号 $B$ をラベルとする重い純な節点 $u$ を固定し，その部分木の重い純な節点の個数を $H_u$ とする．
これらの節点はその部分木の中で祖先について閉じている．
$u$ の生成語のすべての文字は，重い純な節点が書くか，その重くない子の生成語に属する．
各重い純な節点は高々 $\theta$ 文字を書き，そのような子を高々 $r$ 個もち，その各々の総生成語の長さは高々 $\theta+f-1$ である；
したがって $\ell_B\le(r+1)(\theta+f-1)H_u$ である．
これらの節点を，部分木の中での重い純な真の祖先の個数で分類する．
重い純な節点の鎖は (P3) から互いに異なる非終端記号をもつので高々 $v$ 個の節点をもち，
したがってこの個数は $v$ 未満で，分類は高々 $v$ 個である；比較可能な二つの節点は異なる個数をもつので，各分類は反鎖である．
一つの分類の中で，頻出のラベルをもつ節点は長さ少なくとも $\alpha$ で，$u$ の生成語の中に互いに素な生成語をもつので，
高々 $(1+\alpha)\theta/\alpha=\theta+\theta/\alpha<\theta+1$ 個，すなわち高々 $\theta$ 個であり，高々 $v$ 個の分類全体で高々 $v\theta$ 個である；他のラベルをもつ節点はラベルごとに高々 $\theta$ 個なので，$T$ 全体で高々 $v\theta$ 個である．
したがって $H_u\le2v\theta\le v(2\theta+1)$ であり，$\alpha\le g_B\le\ell_B\le\alpha_*$ である．

(P6) \emph{数え上げ．}
混合でない節点で $\ell_u>T_*:=(\alpha_*+1)\theta\le(3f+1)(r+1)v\theta^3$ をみたすものを\emph{大きい}という．
(P4)--(P5) から頻出の非終端記号は大きい節点のラベルにならないので，大きい節点は高々 $v\theta$ 個であり，
それらは混合でない節点の中で祖先について閉じている．
非削除性から節点の生成語は親の生成語の中に現れるので，混合の節点は祖先について閉じている．
$W$ の文字 $\mathtt a$ が，混合でも大きくもない節点 $w$ によって書かれるとする．
二つの閉包性から，混合でも大きくもない $w$ の祖先（$w$ を含む）は，根から $w$ への道の最後の区間をなす；その最上の節点を $z$ とする．
このとき $z$ は $T$ の根か，混合または大きい節点の子であり，$z$ の部分木は混合の節点も大きい節点も含まないので $\ell_z\le T_*$ である．
したがって $W$ のすべての文字 $\mathtt a$ は，混合の節点が書くか，大きい節点が書くか，そのような $z$ の生成語の中にある；
節点 $z$ は高々 $1+rM+rv\theta$ 個である．したがって
\[
\begin{aligned}
 p+K&\le\theta M+(rM+rv\theta+1)T_*+\theta\cdot v\theta\\
 &\le3v\theta^2H_\theta+4rv\theta H_\theta\cdot(3f+1)(r+1)v\theta^3
 \le\bigl(4r(3f+1)(r+1)+3\bigr)v^2\theta^4H_\theta
\end{aligned}
\]
である．ここで $M\le2v\theta H_\theta$ と $rM+rv\theta+1\le4rv\theta H_\theta$ を使った．
$(2f+1)(2f+2)\le D_f$ なので，どちらの場合も，また $K\le\theta+2(f-1)$ のときも，$\tau_X\le C_{f,r}v^2\theta^4H_\theta$ である．
\end{proof}

\begin{proof}[定理~\ref{thm:two-letter-content} の証明]
$L$ が Parikh 飽和ならば補題~\ref{lem:app-two-saturated} が適用できる．
そうでなければ命題~\ref{prop:app-two-nonsaturated} は $\tau(G,c)\le t_{\mathrm{typ}}(G,c)\le D_f(r+1)v\theta$ を与える．
どちらの場合も $\tau(G,c)\le C_{f,r}v^2\theta^4H_\theta$ である．
$c$ の像は三つの元をもつので $|\mathsf S^\sharp|=4$ であり，$N_{\mathrm{typ}}\le C_c(G)\le4^{fr}\|G\|$ である．
系~\ref{cor:distribution-representatives} から，
$\|K\|_+\le N_{\mathrm{typ}}(1+2\theta+r\tau(G,c))+\mathbf 1_{\{\eps\in L\}}$ をみたす共通の特徴標本が存在する．
$v\le\|G\|$ とすると，これは $1+2\theta\le3\|G\|^2\theta^4H_\theta$ と $\tau\le C_{f,r}\|G\|^2\theta^4H_\theta$ を使って最初の評価を，
$1+2\theta\le3\|G\|\theta$ と $\tau\le D_f(r+1)\|G\|\theta$ を使って二番目の評価を与える．
\end{proof}
